%% Generated by tools/build_guide.py from UNIT-REGISTRY.md.
%% Do not edit by hand: edit the source Markdown or the registry.
\chapter{Language Basics}
\label{ch:language-basics}

\guidechapterintro{This chapter introduces the core Croqtile DSL you use to write programs: how a \texttt{\_\_co\_\_} function is declared, how shaped data is described, how iteration and tile selection are expressed, how tiles move, and how concurrent regions coordinate. It stays with short, readable programs and asks only what each construct means.}

\section{The Two Parts of a Croqtile Program}
\label{sec:2-1}

\guideslotintro{A \file{.co} file reads like one C++ program, but the compiler does not treat every function in it the same way. The \texttt{\_\_co\_\_} prefix marks the functions Croqtile rewrites; the rest reaches the downstream compiler as ordinary C++. That split is the first thing to look for when reading an unfamiliar kernel.}

At its simplest a program has two parts. A \texttt{\_\_co\_\_} function describes computation on shaped tensors, and the compiler translates that function into target code. The host program is standard C++ that prepares the data, calls the function and checks the result. Both live in a single \file{.co} file, and the \texttt{\_\_co\_\_} prefix is all the compiler needs to tell them apart.


\subsection*{Host, Device and Tileflow}

Both of those parts are present here, but the ordinary C++ side is worth splitting once more: a \file{.co} file may also carry a \texttt{\_\_device\_\_} helper. The listing is written from the inside out --- a device helper, the \texttt{\_\_co\_\_} function that calls it, and the host \texttt{main()} that calls the \texttt{\_\_co\_\_} function.

\begin{lstlisting}[language=croqtile,caption={The three regions of one source file}]
__device__ void device_function(...) {
  // High-performance device kernel implementation
}
// Tileflow program: orchestrating data movement
__co__ void croqtile_function(...) {
  // ... croqtile code ...
  device_function(...);
  // ...
}
// Host program: typically runs on CPU
void main() {
  // ... prepare data ...
  croqtile_function(...);
  // ...
}
\end{lstlisting}

Execution starts at the bottom of the listing, in the host program. Written in standard C++ and run on the CPU, it is the entry point of the module and the caller of the Tileflow program --- the \texttt{\_\_co\_\_} functions --- and it manages the overall workflow of the heterogeneous application. A typical host body prepares the input data, invokes a Croqtile function, handles its return value and advances to the next step.

The Tileflow program called from there is the part Croqtile adds. Built from \texttt{\_\_co\_\_} functions, it is the core of a \file{.co} program: it orchestrates data movement between host and device and between storage levels within a single device. A typical Tileflow body moves data into a buffer at the appropriate storage level, calls a device function to compute on it, and moves the result back to the host. It is the only region the compiler translates.

Device functions sit at the other end of that call chain. A \texttt{\_\_device\_\_} function defines computation-intensive operations that run on the target device; the prefix marks it as device code. Like the host program, it passes through Croqtile compilation unchanged. The compiler leaves both sides of ordinary C++ alone; only the \texttt{\_\_co\_\_} function between them is rewritten.



\section{Reading a Kernel Line by Line}
\label{sec:2-2}

\guideslotintro{An element-wise kernel is four lines long: a signature, an output allocation, one parallel statement and a return. Reading them one at a time shows why the shape in the signature matters --- it is what lets the compiler check every indexing operation that follows.}

\subsection*{The Croqtile Function}

Start with the signature, which declares more than a C++ function prototype would: every parameter and the result carry a shape.

\begin{lstlisting}[language=croqtile,caption={A kernel signature}]
__co__ s32 [4, 8] ele_add(s32 [4, 8] lhs, s32 [4, 8] rhs) {
\end{lstlisting}

The \texttt{\_\_co\_\_} prefix marks this as a Croqtile function. Unlike regular C++ functions, the signature carries full shape information: \texttt{s32 [4, 8]} means "a 2D tensor of shape 4 \(\times\) 8 with element type \texttt{s32} (signed 32-bit integer)". Both parameters and the return type follow this convention --- \texttt{type [shape] name}. The compiler uses these shapes to verify that every indexing operation is in-bounds at compile time.

Croqtile supports these element types: \texttt{s32} (signed 32-bit int), \texttt{f16} (half-precision float), \texttt{bf16} (brain float), \texttt{f32} (single-precision float), and \texttt{f8\_e4m3} / \texttt{f8\_e5m2} (8-bit floats). In the examples that follow, \texttt{s32} is the simplest to reason about.

\begin{lstlisting}[language=croqtile,caption={Allocating the output}]
s32 [lhs.span] output;
\end{lstlisting}

This allocates the output tensor. The expression \texttt{lhs.span} copies the full shape from \texttt{lhs}, so \texttt{output} automatically has shape \texttt{[4, 8]}. If you later change the input shape, the output adjusts --- a pattern that makes Croqtile functions easier to generalize.

\begin{lstlisting}[language=croqtile,caption={The parallel body}]
parallel {i, j} by [4, 8]
  output.at(i, j) = lhs.at(i, j) + rhs.at(i, j);
\end{lstlisting}

\texttt{parallel \{i, j\} by [4, 8]} creates 4 \(\times\) 8 = 32 parallel instances, one per element position. Each instance runs the body with its own \texttt{(i, j)} pair. The \texttt{.at(i, j)} accessor indexes into a tensor at position \texttt{(i, j)}, and the addition is direct --- one element, one thread.

What \texttt{parallel} leaves unsaid is how those instances map onto the machine: without an annotation the compiler decides. That mapping is deliberately absent from the language here --- the statement describes independent work, not the resources that run it.

\begin{lstlisting}[language=croqtile,caption={Returning the result}]
return output;
\end{lstlisting}

A \texttt{\_\_co\_\_} function returns its result tensor. The return type in the signature (\texttt{s32 [4, 8]}) must match what you actually return. The caller on the host side receives a \texttt{croq:\allowbreak{}:\allowbreak{}spanned\_\allowbreak{}data} --- an owning buffer with shape metadata that you can index into with \texttt{[i][j]}.


\subsection*{\texttt{\_\allowbreak{}\_\allowbreak{}co\_\allowbreak{}\_\allowbreak{} auto} return type}

Use \co@auto@ when the result shape is easiest to describe inside the function. In \co@__co__ auto matmul(...)@, the compiler infers the return type from the buffer named by \co@return output@. The result still has an element type and a shape; \co@auto@ saves repeating that information in the signature. An explicit return type is useful when you want the signature itself to state the required output shape.

The return convention is the same for element-wise arithmetic, scalar matrix multiplication, and MMA kernels that return a buffer. It does not depend on how the function computes the elements.



\section{The Host-Side API: Inputs and Outputs}
\label{sec:2-3}

\guideslotintro{The host creates shaped buffers, passes them to a \texttt{\_\_co\_\_} function and receives a result back. Each step uses a small API, and the shape attached to a buffer is what makes the call type-safe.}

\subsection*{The Host Program}

The host side is plain C++ with two Croqtile API calls: one creates a shaped buffer from raw storage, the other hands a view of it to the kernel.

\begin{lstlisting}[language=croqtile,caption={Creating the shaped inputs}]
auto lhs = croq::make_spandata<croq::s32>(4, 8);
auto rhs = croq::make_spandata<croq::s32>(4, 8);
lhs.fill_random(-10, 10);
rhs.fill_random(-10, 10);
\end{lstlisting}

\texttt{croq:\allowbreak{}:\allowbreak{}make\_\allowbreak{}spandata\textless{}T\textgreater{}(dims.\allowbreak{}.\allowbreak{}.\allowbreak{})} creates an owning tensor buffer on the host. You pass the element type as a template parameter and the dimensions as arguments. \texttt{fill\_random} populates it with values in the given range.

\begin{lstlisting}[language=croqtile,caption={Calling the kernel}]
auto res = ele_add(lhs.view(), rhs.view());
\end{lstlisting}

Calling a \texttt{\_\_co\_\_} function from C++ looks like a normal function call. The \texttt{.view()} method produces a non-owning \texttt{spanned\_view} from an owning \texttt{spanned\_data} --- this is how you pass host tensors into a Croqtile function without transferring ownership. The return value \texttt{res} is a \texttt{croq:\allowbreak{}:\allowbreak{}spanned\_\allowbreak{}data} that owns its buffer.

The rest of \texttt{main} is ordinary verification: loop over every element and assert equality. The host program is deliberately boring --- the interesting work happens in the \texttt{\_\_co\_\_} function.


\subsection*{Host Code}

The host side uses the same \texttt{make\_spandata} / \texttt{.view()} pattern from earlier chapters:

\begin{lstlisting}[language=croqtile,caption={Host Code}]
int main() {
  auto lhs = croq::make_spandata<croq::s32>(128, 256);
  auto rhs = croq::make_spandata<croq::s32>(256, 256);
  lhs.fill_random(-10, 10);
  rhs.fill_random(-10, 10);
  auto res = matmul(lhs.view(), rhs.view());
  for (size_t i = 0; i < res.shape()[0]; ++i)
    for (size_t j = 0; j < res.shape()[1]; ++j) {
      int ref = 0;
      for (size_t k = 0; k < lhs.shape()[1]; ++k)
        ref += lhs[i][k] * rhs[k][j];
      croq::croq_assert(ref == res[i][j], "values are not equal.");
    }
  std::cout << "Test Passed\n" << std::endl;
}
\end{lstlisting}

Nothing new here. The host program stays boring so the Croqtile function stays the star.


\subsection*{Parameter Passing between Croqtile and C++ Function}

Passing arguments between Croqtile and C++ requires the \texttt{choreo.h} header. Combine a raw C++ pointer with its dimensions to construct Croqtile spanned data, as shown below.

\begin{lstlisting}[language=croqtile,caption={Parameter Passing between Croqtile and C Function}]
#include "choreo.h"
void bar(const float* data, unsigned size) {}
__co__ void foo(f32 mdspan<2> d) {
  parallel p by 6 {
    call bar(d, |d|);     // Call the C++ function
  }
}
void foobar(float* a) {
  foo(croq::make_spanview<2>(a, {1, 2}));
}
\end{lstlisting}

The \texttt{make\_spanview} utility template wraps the data without copying it. In the Croqtile function \texttt{foo}, the corresponding parameter \texttt{d} has type \texttt{f32 mdspan\textless{}2\textgreater{}}. An implicit conversion allows \texttt{d} to be passed as the first argument to the C++ function \texttt{bar}; \texttt{|d|} supplies its total element count.

Scalar values can also pass between C++ and Croqtile functions. An \texttt{ituple}, however, is used only inside a Croqtile function.








\subsection*{Input Parameters}

A Croqtile function accepts only two kinds of input parameter, the \emph{spanned data} and the \emph{integer}.

Similar to any C++ program, the parameter may either be named or unnamed:

\begin{lstlisting}[language=croqtile,caption={Input Parameters}]
__co__ void foo(f32 [7] a, int) {...}
\end{lstlisting}

In this example, the second parameter is unnamed, which makes it unable to be reference in the Croqtile function. However, this is still the valid code.


\subsection*{C++ Host: Data Argument with Shape}

One obvious gap between C/C++ pointers/arrays and Croqtile \emph{spanned data} is that \emph{spanned data} is shaped. Therefore, to call a Croqtile function, it is necessary to convert pointers or arrays.

One direct method is to use the Croqtile API provided in \texttt{croqtile.h}. Though programs does not necessary to explicitly include \texttt{croqtile.h} in Croqtile compilation environment, it is included for every Croqtile programs.

\begin{lstlisting}[language=croqtile,caption={C Host: Data Argument with Shape}]
__co__ void foo(f32 [7, 16] input, f32 [7, 16] output) {...}
void entry(float * input, float* output) {
  // foo(input, output); <-- result in error since missing shape
  foo(croqtile::make_spanned<2>(input, {7, 16}),
      croqtile::make_spanned<2>(output, {7, 16}));
}
\end{lstlisting}

In the above example, you make use of \texttt{croqtile:\allowbreak{}:\allowbreak{}make\_\allowbreak{}spanned} function template, providing the rank as template parameter, and typed pointer together with shape list as the function parameters. In this way, a shape is attached with the pointer, which makes a valid Croqtile function parameter. As you may figure out from \texttt{croqtile.h}, it results in \texttt{croqtile:\allowbreak{}:\allowbreak{}spanned\_\allowbreak{}view} object, which simply wraps the C/C++ pointer/array and the array-based list. To make it easy to understand, you list simplified definition is like:

\begin{lstlisting}[language={},caption={C Host: Data Argument with Shape}]
template <typename T, size_t Rank>
class spanned_view {
  T* ptr;                 // typed pointer to the data
  size_t dims[Rank];      // ranked dimensions
};
\end{lstlisting}

It is also possible to have variable as dimensions to construct the shaped data, like following:

\begin{lstlisting}[language=croqtile,caption={C Host: Data Argument with Shape}]
__co__ void foo(f32 [7, 16] input) {...}
void entry(float * input, float* output, int M, int N) {
  foo(croqtile::make_spanned<2>(input, {M, N})),
}
\end{lstlisting}

In this example, the \texttt{input} shape is runtime variable. But do not worry if wrong parameter is provided. Croqtile will check if \texttt{M == 7} and \texttt{N == 16} when entering into the tileflow program. \textbf{If not, the program will terminate immediately} since all the assumptions about shapes in the Croqtile function does not hold.

One important notice for the \emph{spanned data} inputs is that they are \textbf{references to the input data/buffers}. That implies no data copy happens at invoking Croqtile function, though somebody may think from the language perspective it look like copy-semantics.





\subsection*{Return Type and Type Deduction}

Similar to the input, \emph{spanned data} and \emph{integer} are valid output types. The \emph{void} type is also allowed if nothing is returned.

In addition to the explicit return type annotation, Croqtile support type inference on returns. For example:

\begin{lstlisting}[language=croqtile,caption={Return Type and Type Deduction}]
__co__ auto foo() {
  f32 [7, 16] result;
  // ...
  return result;
}
\end{lstlisting}

In this code, the return type is annotated as \texttt{auto}, which is the same keyword as C++. Croqtile is capable of inferring return type as \texttt{f32 [7, 16]}. Therefore, it is recommended to use \texttt{auto} whenever possible.


\subsection*{C++ Host: Shape Return}

The C++ host code that receives the output of Croqtile function actually receives a buffer. Semantics, the return must be copied than reference, or else it returns a reference to an object allocated in the called function. For example:

\begin{lstlisting}[language=croqtile,caption={C Host: Shape Return}]
__co__ auto foo() {
  f32 [7, 16] output;
  // ...
  return output;  // semantically copy
}
void entry() {
  auto result = foo(); // move the 'spanned_data' to caller
  // ...
}
\end{lstlisting}

In this example, the \texttt{output} is semantically copied to the caller. However, in Croqtile's implementation, no copy happens. The returned type is \texttt{croqtile:\allowbreak{}:\allowbreak{}spanned\_\allowbreak{}data}, where its simplified definition is as following:

\begin{lstlisting}[language={},caption={C Host: Shape Return}]
template <typename T, size_t Rank>
class spanned_data {
  std::unique_ptr<T[]> ptr; // move only pointers
  size_t dims[Rank];        // ranked dimensions
};
\end{lstlisting}

In contrary to \texttt{croqtile:\allowbreak{}:\allowbreak{}spanned\_\allowbreak{}view}, \texttt{croqtile:\allowbreak{}:\allowbreak{}spanned\_\allowbreak{}data} owns the buffer it points to. Thus there is \textbf{no copy actually happen when returning from a Croqtile function}. In compiler's term, this is called the \emph{Return Value Optimization}.

You may use member function like \texttt{.rank()}, \texttt{.shape()} to query shape related information of Croqtile return. This allows you able to work without knowing the too much detail about Croqtile function.



\section{Shapes and mdspan}
\label{sec:2-4}

\guideslotintro{A shape is a value a program can name and compute with, not just an annotation on a declaration. \texttt{mdspan} defines one, arithmetic derives another from it, and an anonymous dimension leaves a size for the caller to supply.}

\subsection*{\texttt{span} and \texttt{span(i)}}

\texttt{tensor.span} gives the full shape of a tensor. \texttt{tensor.span(i)} gives one dimension.

\begin{lstlisting}[language=croqtile,caption={span and span(i)}]
s32 [lhs.span(0), rhs.span(1)] output;
\end{lstlisting}

This comes up constantly in data movement code because output shapes often come from different inputs, and moved tiles frequently need to preserve or recompute shape information.

\begin{figure}[htbp]\centering
\includegraphics[width=0.86\linewidth,height=0.42\textheight,keepaspectratio]{f04-fig_span_light.png}
\end{figure}
\begin{center}\small\emph{Use \texttt{span} for the whole shape, \texttt{span(i)} for a single axis.}\end{center}


A shape in Croqtile is a value a program can name and compute with, not just an annotation on a declaration. The examples below define one, derive another from it, and leave a dimension open for the caller to fill in.


\subsection*{Defining Shapes with \texttt{mdspan}}

The keyword \co@mdspan@ stands for multidimensional span. It describes the extents of a shape without allocating its elements. For example, a shape \co@[7, 8]@ has rank 2 and describes 56 positions; choosing \co@f32@ or \co@s32@ for those positions is a separate decision.

You can explicitly define \texttt{mdspan} variables as follows:

\begin{lstlisting}[language=croqtile,caption={Declaring a shape}]
mdspan s0 : [7, 8]; // Defines a 2D shape with dimensions [7, 8]
mdspan<1> s1 : [3]; // Defines a 1D shape with dimensions [3]
\end{lstlisting}

The colon introduces the shape initializer, and square brackets enclose its dimensions. Thus \co@s0@ has seven rows and eight columns, while \co@s1@ has one dimension of size 3. A named shape can be reused in later buffer declarations.

You can specify the \textbf{rank} by placing \texttt{\textless{}rank\textgreater{}} after \texttt{mdspan}. The compiler checks that this rank matches the initialization expression and reports a compilation error if they differ:

\begin{lstlisting}[language=croqtile,caption={A rank mismatch}]
mdspan<3> s2 : [64, 32]; // error: the rank of mdspan is inconsistent
\end{lstlisting}

Because the initializer identifies a shape, the compiler can infer the \co@mdspan@ type when the keyword is omitted:

\begin{lstlisting}[language=croqtile,caption={Omitting the mdspan keyword}]
s3 : [7, 8, 9]
\end{lstlisting}

Here \co@s3@ has rank 3 and dimensions 7, 8, and 9. Every named shape requires an initializer, so the inferred form is often sufficient.


\subsection*{Deriving \emph{mdspan}s}

New shapes often derive from existing ones. \textbf{Tiling} or \textbf{blocking} divides dimensions, while \textbf{padding} increases selected dimensions by a specified amount.

Shape arithmetic expresses these derivations directly:

\begin{lstlisting}[language=croqtile,caption={The original shape}]
shape : [128, 64]; // initial shape
new-shape0 : shape [(0) / 2, (1) / 4, 1];  // tile and reshape: [1, 64, 16]
new-shape1 : shape [(1) + 2, (0) / 16];    // pad and reshape: [66, 8]
\end{lstlisting}

The first derived shape divides dimension 0 by 2 and dimension 1 by 4, then adds a dimension of size 1. The expression determines the new dimension order. It changes the shape description; it does not move any elements.

In Croqtile, the definition of \texttt{new-shape0} is equivalent to:

\begin{lstlisting}[language=croqtile,caption={A derived shape}]
new-shape0: [shape(0) / 2, shape(1) / 4, 1];
\end{lstlisting}

This form lists the dimensions explicitly inside \texttt{[]}. The \textbf{element-of} operation \texttt{()} retrieves a dimension from the existing shape. Both forms produce the same result; the earlier form is syntactic sugar for the complete initialization expression.

The second derived shape adds 2 to the original dimension 1 and divides dimension 0 by 16, then places them in that order.

Furthermore, you may use the \texttt{mdspan} as a whole for derivations:

\begin{lstlisting}[language=croqtile,caption={Rebinding a shape}]
shape : [32, 72]
new-shape0 : shape;  // [32, 72]
new-shape1 : shape + 1;  // [33, 73]
new-shape2 : shape / 4;  // [8, 18]
new-shape3 : [shape, 6]; // [32, 72, 6]
\end{lstlisting}

Arithmetic on a whole shape applies to each dimension. Adding 1 to \co@[128, 64]@ therefore gives \co@[129, 65]@; it does not append a dimension. The two expressions below spell out that same calculation:

The expression \texttt{new-\allowbreak{}shape1 :\allowbreak{} shape + 1;}

is equivalent to

The expression \texttt{new-\allowbreak{}shape1 :\allowbreak{} shape [(0) + 1, (1) + 1];}

The derived definition of \texttt{new-shape3} demonstrates that using an \emph{mdspan} variable in an \emph{mdspan initialization expression} results in \textbf{concatenation} behavior. Thus, the declaration

The expression \texttt{new-\allowbreak{}shape3 :\allowbreak{} [shape, 6];}

is equivalent to

The expression \texttt{new-\allowbreak{}shape3 :\allowbreak{} shape [(0), (1), 6];}

The shorter form is syntactic sugar for the complete shape definition.


\subsection*{\emph{mdspan} with \textbf{Anonymous Dimension}}

Croqtile supports shape dimensions with unknown values, similar to many high-level machine learning languages. The simplest method is to use \textbf{Anonymous Dimension} in \texttt{mdspan}:

\begin{lstlisting}[language=croqtile,caption={mdspan with Anonymous Dimension}]
__co__ auto foo(s32 [?, 1, 2] input) { ... }
\end{lstlisting}

Here, the question mark \texttt{?} represents a compile-time unknown value. Any positive integer could be assumed here. However, the value must be provided at runtime, as the following example shows:

\begin{lstlisting}[language={},caption={mdspan with Anonymous Dimension}]
void bar(int * data, int dim0) {
  foo(croqtile::make_spanview<3>(data, {dim0, 1, 2}));
}
\end{lstlisting}

In this example, the first dimension of \texttt{input} is provided by \texttt{foo}'s \emph{spanned} input, where the value is retrieved from \texttt{bar}'s caller at execution. You may assume the value be determined by reading a configuration file, through shape derivation in a machine learning framework, or etc..

The \emph{anonymous dimension} can work properly for this case. However, in certain scenarios, it is not sufficient.


\subsection*{Step Further: \textbf{Symbolic Dimension}}

In Croqtile, \textbf{symbolic dimensions} give names to dimensions whose values are unknown at compile time. For example:

\begin{lstlisting}[language=croqtile,caption={Step Further: Symbolic Dimension}]
__co__ auto Matmul(s32 [M, K] lhs, s32 [K, N] rhs) {
  int tile_m = 32, tile_m = 8, tile_k = 16;
  s32 [M / tile_m, K / tile_k] tiled_lhs;
  s32 [K / tile_k, N / tile_n] tiled_rhs;
  // ...
}
\end{lstlisting}

Each input dimension has a symbolic name: \texttt{M}, \texttt{K}, or \texttt{N}. Unlike anonymous dimensions, these names make relationships between shapes explicit, improving readability and maintenance. They also give the compiler more information for consistency checks.



\section{Integers, I-Tuples, and Bounded Values}
\label{sec:2-5}

\guideslotintro{Integers control loops and index arrays; an \texttt{ituple} groups several of them into one index. Both are declared with an initializer, both are immutable, and both are checked against the shape they are used with.}

\subsection*{Scalar Types}

Scalar types in Croqtile include the \emph{Integer Type} and \emph{Boolean Type}. The \emph{Integer Type} is similar to C++ type 'int' or 'int32\_t'. It is an signed value which takes 32-bits, ranging from $-2^{31}$ to $2^{31}-1$. The below code illustrates its usage for defining the data and function declaration.

\begin{lstlisting}[language=croqtile,caption={Scalar Types}]
int a;
__co__ int foo(int b);
\end{lstlisting}

Operations like integer arithmetics, shift and logical operation are all supported. And the syntax is identical to C++ builtin operations. In Croqtile, you neither provide equivalence of unsigned scalar integers, nor equivalence of 8-bits, 16-bits, 64-bits scalar integers. This design is simply motivated by the fact: these types are not essential for program control purposes. And normally a 32-bits signed integer can fulfill such work.

The \emph{Boolean Type} in Croqtile is similar to C++ type 'bool'. The operations on top of Boolean, and conversion between Integers, are identical to these of C++.


\subsection*{I-Tuple Types}

An integer tuple is an unordered set of integers. As described, it is normally used as a (subscription) index.

To define an i-tuple, simply enclose its elements within '\{' and '\}' braces. That is,

\begin{lstlisting}[language=croqtile,caption={I-Tuple Types}]
ituple index = {5, 4, 3, 2, 1};  // It defines a tuple of 5 elements
index = {a, b};  // 'a' and 'b' are existing integers
\end{lstlisting}


\subsection*{Operations over \emph{mdspan} and \emph{i-tuple}}

Croqtile supports operations between \texttt{ituple} and \texttt{mdspan} values. The following example applies fixed tiling factors to a shape:

\begin{lstlisting}[language=croqtile,caption={Operations over mdspan and i-tuple}]
sp : [6, 8];
tiling_factor = {3, 2};
spn : sp / tiling_factor;   // spn is defined as [2, 4];
\end{lstlisting}

In Croqtile, mdspan can be defined with such \textbf{Tuple-Span Operations}. The supported operations include:

\begin{itemize}[leftmargin=1.6em]
\item \emph{mdspan} $/$ \emph{i-tuple}
\item \emph{mdspan} $+$ \emph{i-tuple}
\item \emph{mdspan} $\%$ \emph{i-tuple}
\item \emph{mdspan} $*$ \emph{i-tuple}
\item \emph{mdspan} $-$ \emph{i-tuple}
\end{itemize}

Each operation applies to corresponding dimensions. Tuple-span notation collects those calculations into one expression, so the relationship between a shape and its tiling factors remains visible.


In this section, you will learn how to define \emph{integer}s and integer-tuples (\emph{i-tuple}s) in Croqtile code and understand their usage.


\subsection*{Integers for Program Control}

In high-performance computation kernels, integers are typically used for program control rather than computation. In Croqtile \emph{tileflow programs}, integers are used exclusively for loop control and array indexing. Given the limited value ranges in these scenarios, Croqtile has provided a single integer type --- the \emph{32-bit signed integer} --- to simplify the design for this domain-specific use.

In Croqtile, an \textbf{integer}-typed variable is defined similar to C/C++:

\begin{lstlisting}[language=croqtile,caption={Integers for Program Control}]
int a = 1;
\end{lstlisting}

However, since Croqtile can infer the type from the \emph{initialization expression}, you can also define an integer variable in another style:

\begin{lstlisting}[language=croqtile,caption={Integers for Program Control}]
a = 1;
\end{lstlisting}

Similar to C/C++, you can define multiple integers in a single line:

\begin{lstlisting}[language=croqtile,caption={Integers for Program Control}]
int a = 1, b = 2, c = 3;
d = 4, e = 5;
\end{lstlisting}

In Croqtile, the \emph{mdspan}-typed shape can not be used as the function parameter. However, integers can be passed from the host to the Tileflow program, and from the Tileflow program to the device:

\begin{lstlisting}[language=croqtile,caption={Integers for Program Control}]
__co__ void foo(int a) {
  // ...
  call kernel(a, 3);
}
\end{lstlisting}

Unlike C/C++, Croqtile does not allow the declaration of integers without initialization or the reassignment of an integer. For example:

\begin{lstlisting}[language=croqtile,caption={Integers for Program Control}]
int a;  // error: declaration without initialization
int b = 1;
b = 3;  // error: re-assignment
\end{lstlisting}

This enforces that \textbf{integers must be initialized at declaration and cannot be reassigned}. In some contexts, these integers are referred to as "\textbf{immutable}" or "\textbf{constant}".





\subsection*{Defining I-Tuples}

An integer tuple, written \co@ituple@, groups several integer values in a fixed order. Its initializer uses braces:

\begin{lstlisting}[language=croqtile,caption={An i-tuple literal}]
ituple a = {1, 2, 3};
b = {4, 5, 6};  // utilize the type inference
\end{lstlisting}

Since Croqtile can infer types from the \emph{initialization expression}, programmers can often omit the \texttt{ituple} keyword. However, without explicit type annotation, an \emph{ituple} variable definition might look similar to an \emph{mdspan} definition if you are not yet familiar with Croqtile. To distinguish them:

\begin{itemize}[leftmargin=1.6em]
\item The \emph{initialization expression} of an \emph{ituple} follows an assignment operation \texttt{=} (like integers), whereas an \emph{mdspan} is initialized after \texttt{:};
\item The \emph{initialization expression} of an \emph{ituple} is enclosed by \texttt{\{\}}, not \texttt{[]} as for \emph{mdspan}.
\end{itemize}

Similar to \texttt{mdspan}, you may enforce rank check for \emph{ituple}s at compile time:

\begin{lstlisting}[language=croqtile,caption={An inconsistent rank}]
ituple<3> a = {1, 2};  // error: inconsistent rank
\end{lstlisting}


\subsection*{Operations on I-Tuples}

Operations on \emph{i-tuple}s are similar to those on \emph{mdspan}. You can either use the \emph{element-of} operation \texttt{()} to retrieve the element values or use \emph{ituple} as a whole:

\begin{lstlisting}[language=croqtile,caption={Assigning a new tuple}]
a = {3, 4};
b = {a(0), 1, a(1)};  // '()' to retrieve the element value
c = a {(0), (1), 2};  // syntax sugar
d = {a, 5, 6};        // concatenate
e = a + 1;            // addition is applied elementwise
\end{lstlisting}

Similar to \emph{integer} and \emph{mdspan}, an \emph{ituple} variable must be initialized and cannot be reassigned.

In practice, \emph{ituple}s are often used with \emph{mdspan}s. For example:

\begin{lstlisting}[language=croqtile,caption={Indexing an mdspan with an i-tuple}]
shape : [7, 18, 28];
tiling_factors = {1, 2, 4};
tiled_shape : shape / tiling_factors;
padded_shape : shape + {2, 0, 2};
\end{lstlisting}

Here, dividing an \texttt{mdspan} shape by an \texttt{ituple} of tiling factors produces \texttt{tiled\_shape}. Adding an anonymous \texttt{ituple} to the original shape produces \texttt{padded\_shape}. The shape and tuple must have matching ranks; otherwise, the compiler reports an error.

\begin{lstlisting}[language=croqtile,caption={An invalid shape expression}]
shape : [7, 8, 9] + {1, 2}; // error: inconsistent rank
\end{lstlisting}


\paragraph*{In short}

Use an integer for one control value and an \co@ituple@ for an ordered group of values. A shape such as \co@[7, 14]@ describes ranges, while a tuple such as \co@{1, 2}@ can describe a position or a pair of tiling factors. The matching ranks let shape arithmetic apply each factor to the corresponding dimension.

Both integers and ordinary tuples must be initialized and cannot be reassigned. This restriction concerns the names used for control and shape calculations. A buffer element is different: the addition example writes to \co@output.at(i, j)@, and a matrix multiplication updates its accumulator. Keeping these roles distinct avoids treating an immutable index as though it were a mutable accumulator.



\section{Spanned Data and Storage Specifiers}
\label{sec:2-6}

\guideslotintro{Data in Croqtile is \emph{spanned}: a buffer together with its shape. This section shows how spanned data is declared, stored and initialised.}

\subsection*{Fully-Typing}

A buffer declaration needs both an element type and a shape. The shape determines how many positions the buffer contains, while the fundamental type determines what each position stores. The examples below reuse a named shape for half-precision data and use a literal shape for single-precision data.

\begin{lstlisting}[language=croqtile,caption={Complete Types}]
ndims : [20, 15];
f32 [10, 10] d0;
f16 [ndims] d1;
\end{lstlisting}

The basic element types include:

\begin{itemize}[leftmargin=1.6em]
\item Unsigned Integers: \emph{u8/u16/u32}
\item Signed Integers: \emph{s8/s16/s32}
\item Floating-points: \emph{f16/bf16/f32}
\end{itemize}

The types \texttt{s32} and \texttt{int} serve different purposes in Croqtile. The fundamental element type \texttt{s32} must be combined with a shape to declare a fully typed buffer.


\subsection*{The Storage Qualifier}

A tile may be shared by cooperating threads or used privately by one thread. Its element type and shape do not express that distinction; a storage specifier does.

The three storage specifiers used in these examples are:

\begin{itemize}[leftmargin=1.6em]
\item \textbf{global},
\item \textbf{shared},
\item and \textbf{local}.
\end{itemize}

For example:

\begin{lstlisting}[language=croqtile,caption={The Storage Qualifier}]
ndims : [20, 15];
local f32 [10, 10] d0;
shared f16 [ndims] d1;
\end{lstlisting}

Choose the storage specifier to match the code that uses the buffer. Defaults depend on the execution context and target; use an explicit specifier when a transfer or a cooperating group requires a particular storage location.


A Croqtile Tileflow program describes data movement, so buffer declarations are fundamental. Each declaration specifies an element type, a shape, and, when needed, a storage location.


\subsection*{Defining a \emph{Spanned Data}}

Defining a \emph{spanned data} requires both the element type and its shape. For example,

\begin{lstlisting}[language=croqtile,caption={A spanned-data declaration}]
f32 [32, 16] data;
\end{lstlisting}

This defines a buffer with a shape of \texttt{[32, 16]}, where the element type is \texttt{f32}, which is the \emph{IEEE-754} single-precision float. It is also possible to use a named \texttt{mdspan} to define the buffer:

\begin{lstlisting}[language=croqtile,caption={A standalone shape}]
shape : [72, 14];
s32 [shape] data0;              // using the named mdspan `shape`
u32 [shape(0), shape(1)] data1; // same shape as the above
u16 [shape + 6] data2;          // shaped as [78, 20]
s16 [shape, 8] data3;           // shaped as [72, 14, 8]
\end{lstlisting}

This example uses the named \texttt{mdspan} \texttt{shape} to declare the two-dimensional buffer \texttt{data0} with 32-bit integer elements. Place the shape inside \texttt{[]}; element access, arithmetic, and concatenation can also appear in the shape expression.

In Croqtile terminology, the element type is referred to as the \textbf{fundamental type}. Croqtile supports multiple fundamental types, including:

\begin{itemize}[leftmargin=1.6em]
\item \textbf{Unsigned Integers}: \texttt{u8}, \texttt{u16}, \texttt{u32}
\item \textbf{Signed Integers}: \texttt{s8}, \texttt{s16}, \texttt{s32}
\item \textbf{Floating Points}: \texttt{f16}, \texttt{bf16}, \texttt{f32}
\end{itemize}

The prefixes \texttt{u} and \texttt{s} mean unsigned and signed, and the numeric suffix gives the bit width. The type \texttt{f16} is IEEE 754 half precision; \texttt{bf16} is the 16-bit brain floating-point format commonly used in machine learning.

The element type \co@s32@ and the control type \co@int@ have different roles, even though both describe 32-bit integers. A declaration that adds a shape creates a buffer; without one, the same element type declares an ordinary control integer:

\begin{lstlisting}[language=croqtile,caption={Element type without a shape}]
s32 a;    // a control integer, the same as 'int a'
s32 [4] b; // a buffer of four elements
\end{lstlisting}

Both forms compile. The shape decides which one you get, not the type name.


In practical code, some buffer declarations must specify a \textbf{storage specifier}, like so:

\begin{lstlisting}[language=croqtile,caption={Storage Specifier}]
global f32 [32, 7, 2] a;
shared u8 [512, 144] b;
local u8 [72, 1024] c;
\end{lstlisting}

Since Croqtile handles storage in a heterogenous context, a buffer definition without a storage specifier defaults to the storage type of the host program, i.e., CPU memory. Other storage specifiers are defined by the target.

\begin{itemize}[leftmargin=1.6em]
\item \emph{global}: Refers to the device's global storage.
\item \emph{shared}: Refers to the device block's shared storage
\item \emph{local}: Refers to the thread-private storage.
\end{itemize}

Other targets may have different definitions. Furthermore, buffers with different storage types have limitations in their declarations. You will explore this in later chapters.


\subsection*{Initialization}

It is possible to initialize a buffer at the declaration site. For example:

\begin{lstlisting}[language=croqtile,caption={Initialization}]
local s32 [17, 6] b1 {0};       // elements are initialized to 0
shared f32 [128, 16] b2 {3.14f}; // elements are initialized to 3.14f
\end{lstlisting}

The syntax for the \emph{initialization expression} of a \emph{spanned data} is straightforward: it encloses an initial value inside brackets following the variable name. However, its functionality is limited --- it always set all elements to a fixed value.


\subsection*{Declaring the Parameters and Return Values}

Spanned data can appear in the parameter and return declarations of a Croqtile function:

\begin{lstlisting}[language=croqtile,caption={An explicit return shape}]
__co__ f16 [7, 8] foo(f32 [16, 17, 5] input) {...}
\end{lstlisting}

Parameter and return declarations resemble variable declarations, but they do not allow operations on an existing \texttt{mdspan}, storage specifiers, or initializers.

One useful built-in member function for the \emph{spanned} parameters is \texttt{.span}, which provides the associated \emph{mdspan} of the \emph{spanned data}. The following code demonstrates how to use it for a buffer declaration:

\begin{lstlisting}[language=croqtile,caption={A deduced return shape}]
__co__ auto foo(f32 [16, 17, 5] input) {
  f32 [input.span / {4, 1, 5} ] buffer;  // declare a buffer with tiled shape
  // ...
}
\end{lstlisting}


\paragraph*{In short}

This section covered defining and managing data/buffers in Croqtile tileflow programs, including declaration syntax, initialization, storage specifiers, and efficient buffer lifetime management by the compiler. This enables us to step further to the next essential topic: the dynamic shape support.


\subsection*{1. \texttt{global}}

A \co@global@ declaration places the buffer in the target's global storage. It is useful for data that must remain accessible beyond one thread or one cooperating block. Declaring the storage does not initialize its contents or copy data into it:

\begin{lstlisting}[language=croqtile,caption={1. global}]
global f32 [100, 200] matrix;
\end{lstlisting}


\subsection*{2. \texttt{shared}}

Example:

\begin{lstlisting}[language=croqtile,caption={2. shared}]
shared f32 [10, 10] tile_data;
\end{lstlisting}


\subsection*{3. \texttt{local}}

It is particularly useful when each thread requires a small amount of data, such as scalar values or small arrays that do not need to be shared with other threads. Example:

\begin{lstlisting}[language=croqtile,caption={3. local}]
local f16 [5] thread_local_data;
\end{lstlisting}


\subsection*{Default Behavior and Usage}

This default behavior simplifies memory management for users who do not need to explicitly define where their data should be stored, but it also means that accessing such data could incur higher latency, especially when working with large arrays or parallel processing units.

\begin{lstlisting}[language=croqtile,caption={Default Behavior and Usage}]
f32 [100, 100] d0;
\end{lstlisting}

If you want to place data in a different memory space (shared or local), you need to specify the appropriate storage qualifier.

\begin{lstlisting}[language=croqtile,caption={Default Behavior and Usage}]
local f32 [10, 10] d0;     // Small data, local to the thread
shared f16 [ndims] d1;     // Intermediate data, shared between threads in a block
\end{lstlisting}


\subsection*{Example Code: Using Storage Qualifiers}

The following sketch combines the storage declarations:

\begin{lstlisting}[language=croqtile,caption={Example Code: Using Storage Qualifiers}]
ndims : [20, 15];
local f32 [10, 10] d0;
shared f16 [ndims] d1;
global f32 [100, 200] matrix;
// Perform computations on these arrays, utilizing the appropriate storage space
dma.copy d0 => d1;
dma.copy d1 => matrix;
\end{lstlisting}

The declarations illustrate three storage locations. The copy lines sketch the direction of movement between them. To turn this sketch into a working transfer, select source and destination views with compatible shapes: the full arrays declared here have different sizes. A storage specifier chooses a location; it does not make differently shaped buffers interchangeable.



\section{Iteration: parallel-by, with-in, and foreach}
\label{sec:2-7}

\guideslotintro{Iteration in Croqtile declares independence rather than a schedule. This section introduces \texttt{parallel}, \texttt{with-in} and \texttt{foreach} through small programs.}

\subsection*{Croqtile's Two-Layer Parallelism Model}

Croqtile gives you two separate knobs:

\begin{enumerate}[leftmargin=1.6em]
\item \textbf{\texttt{parallel}} --- declares that iterations are independent and \emph{may} run concurrently. This is a logical statement.
\item \textbf{Space specifiers} (\texttt{: block}, \texttt{: thread}, \texttt{: group}, \texttt{: group-4}) --- tell the compiler  each parallel axis maps to. This is a physical statement.
\end{enumerate}

There is also \texttt{foreach}, which is the opposite of \texttt{parallel}: iterations run \textbf{sequentially}, in order. Use \texttt{foreach} when later iterations depend on earlier ones --- an accumulation loop, a K-dimension reduction, a pipeline that loads tile \emph{k} before computing with it.

\begin{lstlisting}[language=croqtile,caption={Croqtile s Two-Layer Parallelism Model}]
parallel {px, py} by [8, 16] : block
  foreach {tile_k} in [16]               // sequential -- K-dimension loop
    parallel {qx, qy} by [16, 16] : thread  // concurrent -- mapped to threads
      foreach k in [16]                      // sequential -- inner accumulation
        output += lhs * rhs;
\end{lstlisting}

\begin{figure}[htbp]\centering
\includegraphics[width=0.86\linewidth,height=0.42\textheight,keepaspectratio]{f05-fig2_logical_vs_physical_light.png}
\end{figure}
\begin{center}\small\emph{Left: the logical nesting the programmer writes (parallel = concurrent, foreach = sequential). Space specifiers bridge the two.}\end{center}

If you omit the space specifier --- just \texttt{parallel p by 8} --- the compiler picks a default mapping. That is what Chapters 1 and 2 did, and it worked fine for simple cases. Explicit specifiers are how you take control when performance matters.

Both \texttt{parallel} and \texttt{foreach} support multi-dimensional indices (\texttt{\{a, b\}} syntax), the compose operator (\texttt{\#}), and the extent operator (\texttt{\#p}). Everything from Chapters 1 and 2 still applies.


\subsection*{\texttt{: block} --- parallel blocks}

\begin{lstlisting}[language=croqtile,caption={: block --- parallel blocks}]
parallel {px, py} by [8, 16] : block
\end{lstlisting}

This declares an 8 by 16 grid, giving 128 blocks. Each block must be able to complete its assigned output tile independently; the example does not provide a communication or synchronization protocol between blocks.

Use \texttt{: block} for the outermost tiling of your output. If your output is \texttt{[128, 256]} and each block handles a \texttt{[16, 16]} tile, you need \texttt{128/\allowbreak{}16 x 256/\allowbreak{}16 = 8 x 16 = 128} blocks.


\subsection*{\texttt{: thread} --- Threads Within a Block}

\begin{lstlisting}[language=croqtile,caption={: thread --- Threads Within a Block}]
parallel {qx, qy} by [16, 16] : thread
\end{lstlisting}

This creates 16 \(\times\) 16 = 256 threads within the enclosing block.

Use \texttt{: thread} for the finest-grain parallelism --- the per-element work within a tile.


\subsection*{Nesting Specifiers}

Real kernels nest these levels. A common pattern:

\begin{lstlisting}[language=croqtile,caption={Nesting Specifiers}]
parallel {px, py} by [8, 16] : block {
  parallel {qx, qy} by [16, 16] : thread {
    // 128 blocks x 256 threads = 32,768 threads total
  }
}
\end{lstlisting}

The outer \texttt{parallel} creates the block grid. The inner \texttt{parallel} creates threads within each block. The braces \texttt{\{px, py\}} introduce multi-dimensional indices --- shorthand for two nested \texttt{parallel} lines.


\subsection*{\texttt{with tile\_\allowbreak{}k in 16}}

Opens a scoped region and binds \texttt{tile\_k} as a tile axis with extent 16. Inside the block, \texttt{tile\_k} is the chunk index for \texttt{chunkat} along K, and \texttt{\#tile\_k} is 16.


\subsection*{\texttt{foreach tile\_\allowbreak{}k(1:\allowbreak{})}: sliced iteration}

The range expression \co@(1:)@ starts iteration at index 1 and keeps the original exclusive upper bound. With 16 tiles, it visits 1 through 15. This is useful after a prologue has loaded tile 0: the loop issues each remaining load once, while computation consumes the previously loaded tile. The final loaded tile is handled after the loop.


\subsection*{Parallel Region: the 'parallel-by' Block}

In systems like CPU, it allows of asynchronous thread to realize the parallel execution. However, in Croqtile, it employs the Single Instruction Multiple Data (SPMD) model as it way to realize parallelization.

The syntax of constructing a parallel region is quite different. It employs the C-style bracket and encloses the code for parallel execution within the 'parallel-by' block.

\begin{lstlisting}[language=croqtile,caption={Parallel Region: the parallel-by Block}]
parallel p by 6 {
  // SPMD code
}
\end{lstlisting}

The above code snippet illustrates the method to create a parallel region with Croqtile keyword 'parallel' and 'by'. Here, it assume there are 6 execution threads. Each of the thread execute the same SPMD code inside but with a different 'p' value. Alternatively, if you are more familiar with sequential C/C++ programming, you may consider 'p' as the iteration variable of a loop with 6 iterations. (But any iteration may go first to be executed!)

Despite parallelism, there is one more implication of 'parallel-by'. In the statement, 'p' is an integer associated with its bound [0, 6). You name 'p' as a \textbf{bounded integer} instead of a simple integer. In some special operations like 'chunkat' (explain later), it requires the \emph{bounded-integer} to work properly since the bound is essential for its computation.


\subsection*{The 'with-in' Block and 'where' Clause}

Similar to 'parallel-by', 'with-in' statement can also bind \emph{i-tuples} to a \emph{mdspan}. The below code shows an example.

\begin{lstlisting}[language=croqtile,caption={The with-in Block and where Clause}]
with index in [10, 10] {
  // index is ituple with 2 elements
}
\end{lstlisting}

Here, 'index' is a \emph{i-tuple} with 2 elements. Sometimes programmers prefer that the 2 elements being named. This is possible by using the below syntax.

\begin{lstlisting}[language=croqtile,caption={The with-in Block and where Clause}]
with {x, y} in [10, 10] {...}
\end{lstlisting}

Or even to name both the \emph{i-tuple} and its elements.

\begin{lstlisting}[language=croqtile,caption={The with-in Block and where Clause}]
with index = {x, y} in [10, 10] {...}
\end{lstlisting}

You name 'index' as a \textbf{bounded ituple} in such scenarios.

While 'with-in' may seem like 'parallel-by', it operates differently. One significant difference is 'with-in' statement does not have implication for parallelism. The code block inside 'with-in' statement is sequentially executed. It does nothing more than creating the \emph{bounded-ituple}.

In addition, programmers could append a 'where' clause. For example,

\begin{lstlisting}[language=croqtile,caption={The with-in Block and where Clause}]
with {m, n} in [M, N], {n_p, k} in [N_P, K] where n_p <-> n {
  // matmul implements with m,n,K. n_p is no long useful.
}
\end{lstlisting}

The code snippet requires 'n' and 'n\_p' to have an identical value in all iterations. Thus inside the 'with-in' block, it is possible to replace'n' whenever 'n\_p' is required, or the opposite. Such a facility is useful to program many AI kernels. Programmers should use operation '\textless{}-\textgreater{}' to establish such relations.


\subsection*{The 'foreach' Block}

Once the \emph{bounded-ituple} is defined by the 'with-in' clause, programmers can loop over the bounded-ituples/bounded-integers. In Croqtile, this is simple.

\begin{lstlisting}[language=croqtile,caption={The foreach Block}]
with x in [10] {
  foreach x {
    // do something with each x
  }
}
\end{lstlisting}


\subsection*{Basic Syntax: \texttt{parallel-by}}

The basic syntax for a \texttt{parallel-by} block is as follows:

\begin{lstlisting}[language=croqtile,caption={Basic Syntax: parallel-by}]
parallel p by 6 {
  // SPMD (Single Program, Multiple Data) code
}
\end{lstlisting}

You name the code inside \texttt{\{\}} as the SPMD-style \textbf{parallel execution block}, and the instances of the code as the \textbf{parallel thread}. In this example:

\begin{itemize}[leftmargin=1.6em]
\item \textbf{\texttt{parallel}}: This keyword initiates the \emph{parallel execution block}.
\item \textbf{\texttt{p}}: It is the \textbf{parallel variable}. It contains an integer value identifying the current \emph{parallel thread}. Each \emph{parallel thread} runs the same code but with a different value of \emph{parallel variable} \texttt{p}.
\item \textbf{\texttt{by 6}}: It indicates that there are totally 6 \emph{parallel thread}s. Thus, the \emph{parallel variable} \texttt{p} ranges from \texttt{0} to \texttt{5} across different threads.
\end{itemize}

In some scenarios, programers may not require an explicit \emph{parallel variable} for a simple parallel construct. It is possible to omit it:

\begin{lstlisting}[language=croqtile,caption={Basic Syntax: parallel-by}]
parallel by 2 { ... }
\end{lstlisting}

This invokes two parallel threads to execute. However, it is not possible for the two threads to work on different data, which is normally required by SPMD programs for a data-parallel processing purpose.


\subsection*{Multiple-Level Parallelism}

In Croqtile, you can define multiple levels of parallelism. Here's an example:

\begin{lstlisting}[language=croqtile,caption={One parallel level}]
parallel p by 2 {
  // parallel-level-0
  parallel q by 12 {
    // parallel-level-1
  }
}
\end{lstlisting}

This code defines two levels of parallelism with values \texttt{2} and \texttt{12}, respectively. This means a total of \texttt{2 x 12 = 24} parallel threads are invoked for the task. Each thread is identified by a unique combination of \texttt{p} and \texttt{q}. In Croqtile, you refer to the level of \texttt{p} as \emph{parallel-level-0} and the level of \texttt{q} as \emph{parallel-level-1}.

One key difference observable by programmers is the \textbf{synchronization cost}. This hierarchical design encourages programmers to place the most divergent code (that requires synchronization) within the \textbf{block} level to enhance overall performance.

In Croqtile syntax, you can define multiple-level parallelism in a single line:

\begin{lstlisting}[language=croqtile,caption={Two parallel levels}]
parallel p by 2, q by 12 { ... }
\end{lstlisting}

This achieves the same parallelization as the previous example. It uses the \textbf{comma-separated} \emph{parallel-by} expression, assuming there is no code in the between two parallel levels.


\subsection*{Sub-Level Parallelism}

Targets like  allow a single parallelism level to be divided into up to three sub-levels. In Croqtile, \texttt{ituple} and \texttt{mdspan} are used for this syntax:

\begin{lstlisting}[language=croqtile,caption={Sub-Level Parallelism}]
parallel {px, py, pz} by [1, 1, 2] {
  // higher synchronization cost
  // ...
  parallel t_index = {qx, qy} by [3, 4] {
    // lower synchronization cost
    // ...
  }
  // ...
}
\end{lstlisting}

In this example, two levels of parallelism are defined. The first level (\emph{level-0}) is subdivided into three sub-levels, annotated by the \texttt{\{px, py, pz\}} \texttt{ituple}. The second level (\emph{level-1}) is subdivided into two sub-levels, annotated by the \texttt{\{qx, qy\}} \texttt{ituple}, or \texttt{t\_index} as a whole. For each sub-level, the parallelization count corresponds to the respective value inside the \texttt{mdspan}.

The first element inside the \texttt{ituple} represents the \textbf{least significant parallel variable}. Croqtile adopts this ordering to align with  developers' conventions, rather than for internal consistency.


In this section, you will learn about constructing \emph{bounded variable}s using the \texttt{with-in} statement and creating loops with the \texttt{foreach} statement.


\subsection*{Loops in Croqtile}

In C++, loops are typically constructed using \texttt{for}, \texttt{while}, and \texttt{do-while} blocks. These loops include conditions for termination and usually involve variable manipulations related to loop control. Programmers have the freedom to define exit conditions, loop stepping, and other aspects, allowing for flexible loop structures.

Croqtile uses bounded ranges for these loops. The range is defined once and then visited by \co@foreach@:

\begin{lstlisting}[language=croqtile,caption={A foreach loop}]
foreach index {
  // loop body
}
\end{lstlisting}

Here, \texttt{index} is a \emph{bounded variable} (introduced in the previous section). If \texttt{index} has an upper bound of \texttt{6}, the above code is equivalent to:

\begin{lstlisting}[language={},caption={The equivalent C++ loop}]
for (int index = 0; index < 6; ++index) {
  // loop body
}
\end{lstlisting}

This form suits tile traversal: the loop visits the valid coordinates in a declared range, and the same bounds can be used to derive each tile's shape.

In Croqtile, \emph{bounded variables} defined within a \texttt{parallel-by} statement are immutable and cannot be used in \texttt{foreach} statements. Instead, programmers should use the \texttt{with-in} statement to define the \emph{bounded variables} suitable for loops.





\subsection*{Defining \emph{Bounded Variable}s:}

A \texttt{with-in} statement allows you to define either a \emph{bounded integer} or a \emph{bounded ituple}. Syntactically, it resembles the \texttt{parallel-by} statement but only defines \emph{bounded variable}s without invoking multiple instances for execution. Below are examples:

\begin{lstlisting}[language=croqtile,caption={Defining Bounded Variables:}]
with x in 128 {
  // 'x' is a bounded integer
}
with y in [512] {
  // 'y' is a bounded ituple
}
with index in [10, 10] {
  // 'index' is a bounded ituple
}
\end{lstlisting}

Here, \texttt{x} represents a bounded integer with an upper bound of \texttt{128}. Both \texttt{y} and \texttt{index} represent bounded ituples whose upper bounds are \texttt{512} and \texttt{10, 10} respectively. Similar to the \texttt{parallel-by} statement, you can name the elements of the ituple for clarity or declare both the bounded ituple and the associated bounded integers, as shown below:

\begin{lstlisting}[language=croqtile,caption={Defining Bounded Variables:}]
with {x, y} in [10, 10] {
  // x and y are explicitly named elements of the ituple
}
with index = {x, y} in [10, 10] {
  // 'x', 'y', and 'index' can be used within the block
}
\end{lstlisting}

Like the \texttt{parallel-by} statement, multiple \texttt{with-in} declarations can be combined using a comma-separated syntax, as demonstrated below:

\begin{lstlisting}[language=croqtile,caption={Defining Bounded Variables:}]
with index = {x, y} in [10, 10], idx in [100, 10] { }
\end{lstlisting}

In this way, two bounded ituples are defined. Note that either \texttt{index} or \texttt{idx} can only be referenced within the following \texttt{with-in} block.





The \texttt{with-in} statement defines \emph{bounded variable}s, which can then be iterated over using the \texttt{foreach} statement. The basic syntax is as follows:

\begin{lstlisting}[language=croqtile,caption={Indexing one shape}]
with index in [6] {
  foreach index {
    // Perform operations with each index
  }
}
\end{lstlisting}

In this example, the \texttt{foreach} block iterates 6 times, with the value of the \textbf{iteration variable} \texttt{index} ranging from \texttt{0} to \texttt{5} incrementally.

It is also possible to iterate over \emph{bounded ituple}s. For example:

\begin{lstlisting}[language=croqtile,caption={A tuple of indices}]
with index = {x, y} in [6, 17] {
  foreach index { }
}
\end{lstlisting}

This code defines a \emph{bounded ituple} \texttt{index}. Iterating over \texttt{index} is equivalent to a nested loop structure:

\begin{lstlisting}[language={},caption={The generated C++ loop}]
for (int x = 0; x < 6; ++x)
  for (int y = 0; y < 17; ++y) { }
\end{lstlisting}

The left-to-right order of indices in a bounded tuple determines the outer-to-inner nesting of the loops. Keep this ordering in mind when the order of iteration matters.

The same outer-to-inner ordering applies to a comma-separated list of bounded variables:

\begin{lstlisting}[language=croqtile,caption={Several bounded groups}]
with index = {x, y} in [6, 17], iv in [128] {
  foreach iv, index { }
}
\end{lstlisting}

In this code, two bounded variables are defined. The following \texttt{foreach} block is equivalent to a multi-level nested loop:

\begin{lstlisting}[language={},caption={The generated inner loop}]
for (int iv = 0; iv < 128; ++iv)
  for (int x = 0; x < 6; ++x)
    for (int y = 0; y < 17; ++y) { }
\end{lstlisting}


\subsection*{Syntactic Sugar}

\begin{lstlisting}[language=croqtile,caption={A foreach loop}]
foreach x in 128 { }
foreach idx in [10, 20] { }
foreach {y, z} in [8, 16] { }
\end{lstlisting}

The code above is equivalent to the following code:

\begin{lstlisting}[language=croqtile,caption={The with-in form it expands to}]
with x in 128 {
  foreach x { }
}
with idx in [10, 20] {
  foreach idx { }
}
with {y, z} in [8, 16] {
  foreach y, z { }
}
\end{lstlisting}


\subsection*{Deriving the Loop From a Bounded Integer}

In certain scenarios, such as pipelining data movement, it may be necessary to modify loop iterations. In Croqtile, this can be achieved by deriving a loop from a \emph{bounded integer} within the \texttt{foreach} statement. For example:

\begin{lstlisting}[language=croqtile,caption={A tuple of bounded variables}]
with {x, y} in [6, 17] {
  foreach x, y(1::) { }
}
\end{lstlisting}

In this case, the \textbf{Range Expression} \texttt{y(1::)} is used to derive the loop. This results in equivalent C/C++ code:

\begin{lstlisting}[language={},caption={The derived C++ loop}]
for (int x = 0; x < 6; ++x)
  for (int y = 1; y < 17; ++y) { }
\end{lstlisting}

The \texttt{y} loop starts at \texttt{1}. A range expression consists of a bounded variable followed by three colon-separated integer fields in parentheses:

\begin{lstlisting}[language=croqtile,caption={The bounded-variable syntax}]
  bounded-variable(lower-offset:upper-offset:stride)
\end{lstlisting}

This derives a loop from the \texttt{bounded-\allowbreak{}variable}, where the \texttt{bounded-\allowbreak{}variable} serves as the \emph{iteration variable} of the loop. Specifically:

\begin{itemize}[leftmargin=1.6em]
\item The initial value of the \emph{iteration variable} is the \texttt{lower-offset} plus the lower bound of the \texttt{bounded-\allowbreak{}variable}. Since all bounded integers have a lower bound of \texttt{0} in Croqtile, the \texttt{lower-offset} sets the initial value of the loop.
\item The loop terminates when the \emph{iteration variable} is equal to or greater than the upper bound of the \texttt{bounded-\allowbreak{}variable} plus \texttt{upper-offset}. Negative values are typically used for \texttt{upper-offset}.
\item The \emph{iteration variable} increments by \texttt{stride} at the end of each iteration.
\end{itemize}

Thus, \texttt{y(1:-1:2)} results in a loop like \texttt{for (y = 0 + 1; y \textless{} 17 -\allowbreak{} 1; y += 2)} in the example above. If any field of the \emph{range expression} is not specified, it results in default values: \texttt{0} for \texttt{lower-offset} and \texttt{upper-offset}, and \texttt{1} for \texttt{stride}.

Note that \emph{range expression}s only apply to the \emph{bounded variable}s. \emph{Range expression} over the \emph{bounded ituple} triggers an error at compile time.


\paragraph*{In short}

Bounded variables and the loops derived from them are what data movement code is built from. You introduced the syntax and behavior of these constructs, including how to apply range operations to modify loop iterations and the importance of understanding the current and upper-bound values of bounded variables.

The loop deriving part is important for implementing multi-buffering data movement, which is essential for building high-performance kernels and will be introduced in optimization chapters later.



\section{Tile Selection: chunkat, subspan, and view/from}
\label{sec:2-8}

\guideslotintro{A tile is a slice of a larger buffer. This section shows the ways to select one --- \texttt{chunkat}, \texttt{subspan}, \texttt{view().from()} --- and the compose operator \texttt{\#}.}

\subsection*{\texttt{chunkat(...)}: equal partitioning}

\begin{lstlisting}[language=croqtile,caption={chunkat(...): equal partitioning}]
lhs.chunkat(tr, tc)
\end{lstlisting}

\texttt{chunkat} divides each dimension into equal chunks based on the surrounding parallel structure. If \texttt{lhs} has shape \texttt{[64, 128]} and the kernel uses \texttt{parallel \{tr, tc\} by [4, 8]}, then \texttt{lhs.\allowbreak{}chunkat(tr, tc)} selects a \texttt{[16, 16]} tile.

That makes \texttt{chunkat} the natural tool when your tiling follows the parallel decomposition directly.

\begin{figure}[htbp]\centering
\includegraphics[width=0.86\linewidth,height=0.42\textheight,keepaspectratio]{f04-fig3_chunkat_light.png}
\end{figure}
\begin{center}\small\emph{\texttt{chunkat(1, 3)} selects one equal tile from a regular tiling grid.}\end{center}


\subsection*{\texttt{subspan(.\allowbreak{}.\allowbreak{}.\allowbreak{}).\allowbreak{}at(.\allowbreak{}.\allowbreak{}.\allowbreak{})}: explicit tile extents on a tile grid}

\begin{lstlisting}[language=croqtile,caption={subspan(...).at(...): explicit tile extents on a tile grid}]
lhs.subspan(64, 64).at(block_m, iv_k)
\end{lstlisting}

The \co@subspan@ expression supplies the tile extents explicitly, and \co@.at(...)@ supplies tile coordinates. For a \co@[16, 16]@ tile, coordinate \co@(2, 3)@ begins at element \co@(32, 48)@. This form is useful when the tile size is a fixed algorithmic choice rather than derived from the number of parallel instances.


\subsection*{\texttt{view(.\allowbreak{}.\allowbreak{}.\allowbreak{}).\allowbreak{}from(.\allowbreak{}.\allowbreak{}.\allowbreak{})}: explicit tile extents with explicit element offsets}

\begin{lstlisting}[language=croqtile,caption={view(...).from(...): explicit tile extents with explicit element offsets}]
lhs.view(TILE_M, TILE_K).from(offset_m, offset_k)
\end{lstlisting}

\texttt{view(.\allowbreak{}.\allowbreak{}.\allowbreak{}).\allowbreak{}from(.\allowbreak{}.\allowbreak{}.\allowbreak{})} is the most general selection form introduced in this part of the guide. \texttt{view(...)} states the shape of the region you want, and \texttt{.from(...)} gives its starting coordinate in element space.

This is the right tool when the tile origin is already available as element offsets rather than tile indices, or when the region is not most naturally described as an equal chunk. It is especially useful for dynamic and irregular access patterns where \texttt{chunkat(...)} would be awkward or impossible.

You can think of the three forms like this:

\begin{itemize}[leftmargin=1.6em]
\item Use \texttt{chunkat(...)} when the surrounding parallel decomposition already defines an equal tiling.
\item Use \texttt{subspan(.\allowbreak{}.\allowbreak{}.\allowbreak{}).\allowbreak{}at(.\allowbreak{}.\allowbreak{}.\allowbreak{})} when you want an explicit tile shape and tile coordinates.
\item Use \texttt{view(.\allowbreak{}.\allowbreak{}.\allowbreak{}).\allowbreak{}from(.\allowbreak{}.\allowbreak{}.\allowbreak{})} when you want an explicit tile shape and element-space starting offsets.
\end{itemize}


\subsection*{Compose: \texttt{outer \# inner}}

\begin{lstlisting}[language=croqtile,caption={Compose: outer inner}]
output.at(tr # i, tc # j)
\end{lstlisting}

This combines a tile coordinate with an in-tile coordinate to produce a coordinate in the full tensor.

Read \texttt{tr \# i} as "row \texttt{i} inside tile \texttt{tr}."

\begin{figure}[htbp]\centering
\includegraphics[width=0.86\linewidth,height=0.42\textheight,keepaspectratio]{f04-fig_compose_light.png}
\end{figure}
\begin{center}\small\emph{Compose a tile index with an in-tile offset to recover a global coordinate.}\end{center}


\subsection*{Extent: \texttt{\#name}}

\begin{lstlisting}[language=croqtile,caption={Extent: name}]
foreach i in [128 / #tile]
\end{lstlisting}

The notation \texttt{\#tile} means the extent of the \texttt{tile} axis: how many tiles exist along that dimension. If \texttt{parallel tile by 8}, then \texttt{\#tile} is \texttt{8}.

Figure~\ref{fig:f04-fig_extent} shows both uses of the \texttt{\#} prefix.

\begin{figure}[htbp]\centering
\includegraphics[width=\linewidth,height=0.56\textheight,keepaspectratio]{f04-fig_extent_light.png}
\caption{Prefix \#tile gives the axis extent, 8; infix tile \# i composes a tile coordinate and an in-tile offset.}
\label{fig:f04-fig_extent}
\end{figure}


\subsection*{Bounded-ituple/integer and 'chunkat' Operation}

'chunkat' is an operation performed on spanned data. It creates a new \emph{mdspan} over the existing data. Thus in certain systems, it is referred to as 'subview'. However, as 'chunkat' accepts bounded-ituple and bounded-integer as parameters, it is named differently in Croqtile.

\begin{lstlisting}[language=croqtile,caption={Bounded-ituple/integer and chunkat Operation}]
global f32 [6, 10, 100] data;
parallel p by 6 {
  with index = {x, y} in [10, 10] {
    // for every data move, the stride into 'data' is
    //    stride = p*1000 + x * 100 + y * 10
    //
    // for each chunk, the dimensioned size for the movement is {1, 1, 10}
    f = dma.copy data.chunkat(p, index) => local;
  }
}
\end{lstlisting}

The example above shows one typical use of \texttt{chunkat}. Here the spanned data has type \texttt{f32 [6, 10, 100]}. The operation takes two parameters --- the \emph{integer} \texttt{p} and the \emph{ituple} \texttt{index} --- and divides the data into \(6 \times 10 \times 10\) pieces, one for each combination of their ranges. Each piece holds 10 consecutive elements along the least significant dimension and is copied into local storage.


A \emph{chunkat expression} names one tile of a larger buffer in a DMA statement, so the movement code never has to compute an offset by hand.


\subsection*{\emph{Ubound} Operation}

In the previous chapter, you learned that a \emph{bounded variable} defined by a \texttt{parallel-by} or \texttt{with-in} statement is either an \emph{integer} or an \emph{ituple} with one or more bounds. You also learned how to use the \emph{current value} and \emph{upper-bound value} of the \emph{bounded variable} to construct loops in \texttt{foreach} blocks, where the \emph{upper-bound} is implicitly used.

However, in Croqtile, you can use the \textbf{ubound operation} to explicitly retrieve the upper bound of a bounded variable. \emph{Ubound} is a unary operation that uses the prefix symbol \texttt{\#} before a bounded variable. For example:

\begin{lstlisting}[language=croqtile,caption={Ubound Operation}]
with size in [32] {
  local u8 [#size, 16] a; // buffer shape: [32, 16]
  local u8 [size, 16] b; // error in shape: [0, 16]
}
\end{lstlisting}

In this code, the local buffer \texttt{a} is shaped as \texttt{[32, 16]}, since \texttt{\#size} evaluates to the upper-bound of \texttt{size}. However, the declaration of \texttt{b} uses the \emph{current value} of \texttt{size}, resulting in an incorrect shape definition.


\subsection*{Use \texttt{chunkat} Expression to Tile Data}

A transfer can use a whole buffer, but a tiled algorithm usually needs only the portion assigned to the current instance. Put \co@chunkat@ in the source expression to describe that portion directly:

\begin{lstlisting}[language=croqtile,caption={Use chunkat Expression to Tile Data}]
__co__ void foo(f32 [36] input) {
  parallel p by 6 : block {
    dma.copy input.chunkat(p) => shared;
  }
}
\end{lstlisting}

In this example, the built-in member function \texttt{.chunkat} is used to create a tile or block for the source expression of the DMA statement. Each tiled block is called a \textbf{chunk} in Croqtile. The \texttt{.chunkat(...)} data expression can be considered a simpler form of \texttt{.\allowbreak{}chunk(.\allowbreak{}.\allowbreak{}.\allowbreak{}).\allowbreak{}at(.\allowbreak{}.\allowbreak{}.\allowbreak{})}, where the semantics of each part are illustrated in the figure below:

Figure~\ref{fig:f29-chunkat} shows how the indices walk the buffer.

\begin{figure}[htbp]\centering
\includegraphics[width=\linewidth,height=0.56\textheight,keepaspectratio]{f29-chunkat.png}
\caption{\co@chunkat@ divides a buffer into the tiles the parallel instances consume.}
\label{fig:f29-chunkat}
\end{figure}

Considering the expression \texttt{input.\allowbreak{}chunkat(p)}, which is equivalent to \texttt{input.\allowbreak{}chunk(\#\allowbreak{}p).\allowbreak{}at(p)}:

\begin{itemize}[leftmargin=1.6em]
\item \texttt{input} is the \emph{spanned data} to be tiled.
\item \texttt{\#p} indicates the \textbf{tiling factor}. It also specifies the total \textbf{chunk count} for each dimension to process. In this case, the 1D input is tiled by \texttt{6}, resulting in \texttt{6} chunks to be processed.
\item \texttt{.chunk(\#p)} determines the \textbf{chunk shape}. The chunk is a 1D data segment tiled from \texttt{input}, with its size being the data \emph{mdspan} divided by the \emph{upper-bound} of the bounded variable. In this example, the chunk shape is calculated as \texttt{[input.\allowbreak{}span /\allowbreak{} \#\allowbreak{}p] = [6]}.
\item \texttt{.at(p)} determines the \textbf{chunk index}. Here, \texttt{p} ranges from \texttt{[0, 6)}, allowing parallel threads to process different chunks of \texttt{input}. Specifically, parallel thread \texttt{1} processes a chunk with elements \texttt{[6, 12)} from \texttt{input}, while parallel thread \texttt{5} processes a chunk with elements \texttt{[30, 36)}, and so on.
\end{itemize}

In many scenarios, it is unnecessary to use separate \texttt{.chunk} and \texttt{.at} expressions, as \texttt{.chunkat} takes bounded variables as parameters. The upper-bounds of these variables act as the \emph{tiling factors}, and their current values act as the \emph{chunk index}.


\subsection*{Use Multiple Bounded Variables}

The example below uses several bounded variables at once:

\begin{lstlisting}[language=croqtile,caption={Use Multiple Bounded Variables}]
__co__ void foo(f32 [36, 14, 8] input) {
  parallel p by 6 {
    with index in [7, 4] {
      foreach index {
        dma.copy input.chunkat(p, index) => shared;
      }
    }
  }
}
\end{lstlisting}

In this example, there are two bounded variables: \texttt{p} is the \emph{bounded integer} defined by the \texttt{parallel-by} statement, while \texttt{index} is the \emph{bounded ituple} with two integers. The \texttt{chunkat} expression concatenates them to form a \emph{bounded ituple} with upper-bounds \texttt{\{\#\allowbreak{}p, \#\allowbreak{}index\}} and current values \texttt{\{p, index\}}. Thus, the shape of each chunk is calculated as \texttt{[36, 14, 8] /\allowbreak{} \{\#\allowbreak{}p, \#\allowbreak{}index\} = [36, 14, 8] /\allowbreak{} \{6, 7, 4\} = [6, 2, 2]}.

In addition, since \texttt{p} is the bounded variable for parallelization, there are \texttt{6} threads working in parallel, each handling a large chunk of \texttt{[6, 14, 8]}. Within a single parallel thread, it iterates \texttt{7x4} times, each time moving a \texttt{[6, 2, 2]} small chunk to a compiler-allocated \texttt{shared} storage.

With this approach, you can apply tiling to data movement, where each iteration in each parallel thread selects only a part of the data/buffer for transfer. The Croqtile compiler uses information about the \emph{spanned} data and the \emph{bounded variables} to transpile the DMA statement into C++ subscripting, DMA configurations, etc., simplifying the trivial and error-prone work that C++ programmers typically face. This makes learning the Croqtile syntax worthwhile.


\paragraph*{In short}

In this section, you learned about using the \texttt{chunkat} expression in Croqtile for tiling data in DMA statements, which allows dividing data into smaller chunks for efficient processing. You explored how \texttt{chunkat} can be applied to either the source or destination of a DMA statement to manage data movement between memory and processors.

Now, you may try to write a full tileflow program. Wish you good luck!



\section{Data Movement: the DMA Statement and Futures}
\label{sec:2-9}

\guideslotintro{Moving a tile is a statement in the language, not a library call. This section introduces the DMA statement, futures, and synchronous versus asynchronous copies.}

\subsection*{The DMA Statement Family}

All of Croqtile's tile-movement operations use the same basic form:

\begin{lstlisting}[language=croqtile,caption={The general DMA form}]
future = engine.operation.modifiers source => destination;
\end{lstlisting}

The important pieces are:

\begin{table}[htbp]
\centering\small
\begin{tabularx}{\linewidth}{@{}lX@{}}
\toprule
\textbf{Part} & \textbf{Meaning} \\
\midrule
\texttt{future} & Optional handle to the transfer result. Use \texttt{future.data} to access the moved tile and \texttt{future.span} to query its shape. \\
\texttt{operation} & \texttt{copy}, \texttt{transp\textless{}.\allowbreak{}.\allowbreak{}.\allowbreak{}\textgreater{}}, or \texttt{pad\textless{}.\allowbreak{}.\allowbreak{}.\allowbreak{}\textgreater{}}. \\
\texttt{modifiers} & Optional flags such as \texttt{.async}, \texttt{.zfill}, or \texttt{.\allowbreak{}swiz\textless{}N\textgreater{}}. \\
\texttt{source =\textgreater{} destination} & The source view and the destination memory or destination view. \\
\bottomrule
\end{tabularx}
\end{table}

Minimal examples:

\begin{lstlisting}[language=croqtile,caption={A concrete copy}]
f = dma.copy input.chunkat(block) => shared;
g = dma.transp<1, 0> matrix => local;
h = dma.pad<{2, 1}, {3, 2}, {0, 0}, 0> tile => shared;
\end{lstlisting}

The key shift from Section~\ref{sec:1-7} is that these statements move shaped regions, not individual elements.


\subsection*{First Example: Tiled Addition Through shared storage}

\begin{lstlisting}[language=croqtile,caption={First Example: Tiled Addition Through shared storage}]
__co__ s32 [64, 128] tiled_add_2d(s32 [64, 128] lhs, s32 [64, 128] rhs) {
  s32 [lhs.span] output;
  parallel {tr, tc} by [4, 8] {
    lhs_load = dma.copy lhs.chunkat(tr, tc) => shared;
    rhs_load = dma.copy rhs.chunkat(tr, tc) => shared;
    foreach {i, j} in lhs_load.span
      output.at(tr # i, tc # j) =
          lhs_load.data.at(i, j) + rhs_load.data.at(i, j);
  }
  return output;
}
\end{lstlisting}

The arithmetic is still element-wise addition. What changed is the data path:

\begin{enumerate}[leftmargin=1.6em]
\item \texttt{chunkat(tr, tc)} selects one logical tile from each input.
\item \texttt{dma.\allowbreak{}copy .\allowbreak{}.\allowbreak{}.\allowbreak{} =\textgreater{} shared} stages that tile into block-visible fast memory.
\item \texttt{tr \# i} and \texttt{tc \# j} map tile-local coordinates back to global output coordinates.
\end{enumerate}

This is the core Croqtile pattern: select a tile, move it, compute on the moved tile.

\begin{figure}[htbp]\centering
\includegraphics[width=0.86\linewidth,height=0.42\textheight,keepaspectratio]{f04-fig2_tiled_add_light.png}
\end{figure}
\begin{center}\small\emph{Load a tile, compute on the staged data, then write results back using global coordinates.}\end{center}


\subsection*{Futures, \texttt{.data}, and \texttt{.span}}

The result of a DMA statement is a future-like handle:

\begin{lstlisting}[language=croqtile,caption={A future from a copy}]
f = dma.copy input => shared;
\end{lstlisting}

That handle gives you two important things:

\begin{itemize}[leftmargin=1.6em]
\item \texttt{f.data}: the moved tile, as a spanned tensor.
\item \texttt{f.span}: the shape of that tile.
\end{itemize}

Example:

\begin{lstlisting}[language=croqtile,caption={Reading data and span}]
f = dma.copy input => shared;
foreach idx in [f.span]
  output.at(idx) = f.data.at(idx);
\end{lstlisting}

Even synchronous copies are worth naming when you want to pass the moved tile to later code.

\begin{figure}[htbp]\centering
\includegraphics[width=\linewidth,height=0.56\textheight,keepaspectratio]{f04-fig_future_data_light.png}
\caption{The copy result exposes the destination tensor through .data and its shape through .span; asynchronous reads require completion first.}
\end{figure}


\subsection*{Synchronous and Asynchronous Copies}

By default, copies are synchronous:

\begin{lstlisting}[language=croqtile,caption={A synchronous copy}]
f = dma.copy input => shared;
dma.copy f.data => output;
\end{lstlisting}

Add \texttt{.async} when you want the transfer to proceed in the background:

\begin{lstlisting}[language=croqtile,caption={An asynchronous copy}]
f = dma.copy.async input => shared;
wait f;
dma.copy f.data => output;
\end{lstlisting}

The same pattern applies to asynchronous copies:

\begin{lstlisting}[language=croqtile,caption={An asynchronous copy}]
f = dma.copy.async input => shared;
wait f;
dma.copy f.data => output;
\end{lstlisting}

A future identifies the destination before an asynchronous transfer necessarily finishes. Access through \co@f.data@ does not mean the elements are ready. Place \co@wait f@ before the first operation that reads those elements, and keep the source and destination valid while the transfer is in flight.

Later chapters use async copies to overlap load and compute. In this part of the guide, the important part is the contract: \texttt{.async} creates an in-flight transfer, and \texttt{wait} makes it safe to consume.


\subsection*{\texttt{dma.\allowbreak{}transp\textless{}.\allowbreak{}.\allowbreak{}.\allowbreak{}\textgreater{}}}

\begin{lstlisting}[language=croqtile,caption={dma.transp ...}]
fa = dma.transp<1, 0, 2> a => local;
dma.copy fa.data => o;
\end{lstlisting}

This permutes dimensions during the movement itself. You do not first copy and then transpose. The transfer produces data in the transposed layout.


\subsection*{\texttt{dma.\allowbreak{}pad\textless{}.\allowbreak{}.\allowbreak{}.\allowbreak{}\textgreater{}}}

\begin{lstlisting}[language=croqtile,caption={dma.pad ...}]
f = dma.pad<{2, 1}, {3, 2}, {0, 0}, V>.async input => shared;
wait f;
dma.copy f.data => output;
\end{lstlisting}

Padding creates a larger destination tile. The template arguments describe low padding, high padding, internal stride-style padding, and the fill value. Use \texttt{pad} when you want explicit border values, not just out-of-bounds protection.


\subsection*{Tileflow Program: Orchestrating the Data Movement}

The \emph{Tileflow Program} consists of \emph{Croqtile functions}. As described earlier, it manages the movement of data between the host and the target device, ensuring that data is copied correctly across different storage locations.

For convenience, let's revisit the code:

\begin{lstlisting}[language=croqtile,caption={The tileflow program}]
__co__ s32 [6, 17, 128] ele_add(s32 [6, 17, 128] lhs, s32 [6, 17, 128] rhs) {
  s32 [lhs.span] output; // Use same shape as lhs
  parallel p by 6 {
    with index in [17, 4] { // Tiling factors
      foreach index {
        lhs_load = dma.copy lhs.chunkat(p, index) => local;
        rhs_load = dma.copy rhs.chunkat(p, index) => local;
        local s32 [lhs_load.span] l1_out;
        // Call kernel with loaded data
        call kernel(lhs_load.data, rhs_load.data, l1_out, |lhs_load.span|);
        // Store result back to output
        dma.copy l1_out => output.chunkat(p, index);
      }
    }
  }
  return output;
}
\end{lstlisting}

The \texttt{\_\_co\_\_} function accepts \texttt{lhs} and \texttt{rhs}, both with shape \texttt{[6, 17, 128]} and element type \texttt{s32} (signed 32-bit integer). The output has the same shape and element type.

The \texttt{parallel p by 6 \{.\allowbreak{}.\allowbreak{}.\allowbreak{}\}} block indicates that the enclosed code runs in parallel. Specifically, six instances of the code are executed concurrently. This implies a transition in the execution environment from the host to the device.

Inside the \texttt{parallel-by} block, \texttt{with-in} binds \texttt{index} to the bounds \texttt{[17, 4]}. This two-dimensional bounded tuple supplies the indices for the \texttt{foreach} loop.

The \texttt{foreach index \{.\allowbreak{}.\allowbreak{}.\allowbreak{}\}} statement is equivalent to the following C code:

\begin{lstlisting}[language={},caption={The equivalent C++ loop}]
for (int x = 0; x < 17; x++)  // assume 'x' represents the 1st element of 'index'
  for (int y = 0; y < 4; y++) { ... }  // and 'y' represents the 2nd element of 'index'
\end{lstlisting}

Within the \texttt{foreach} block, the \texttt{dma.copy} statement describes how data movement occurs. For example, consider the statement \texttt{lhs\_\allowbreak{}load = dma.\allowbreak{}copy lhs.\allowbreak{}chunkat(p, index) =\textgreater{} local;}:

\begin{itemize}[leftmargin=1.6em]
\item The symbol \texttt{lhs\_load} is called \textbf{future} of the DMA operation, which contains information about the DMA destination.
\item \texttt{dma.copy} invokes a direct DMA data transfer without transforming the shape of the data. The expression on the left-hand side of \texttt{=\textgreater{}} represents the DMA source, while the right-hand side represents the destination.
\item In this case, the destination is specified as a \texttt{local} buffer, which is automatically allocated by the Croqtile compiler.
\item The source expression \texttt{lhs.\allowbreak{}chunkat(p, index)} selects a chunk of \texttt{lhs}. The bounds of \texttt{p} and \texttt{index}, namely \texttt{[6, 17, 4]}, determine the partition. Dividing the input shape \texttt{[6, 17, 128]} by these bounds gives chunk shape \texttt{[1, 1, 32]}. The current values select the chunk: instance \texttt{p = 1} and iteration \texttt{index = \{16, 2\}} select chunk \texttt{(1, 16, 2)}, whose element origin is \texttt{(1, 16, 64)}.
\end{itemize}

The following figure shows how the bounds determine the chunk shape and the indices select a chunk.

Figure~\ref{fig:f14-chunkat_3d} shows how the bounds determine the chunk shape.

\begin{figure}[htbp]\centering
\includegraphics[width=\linewidth,height=0.56\textheight,keepaspectratio]{f14-chunkat_3d_light.png}
\caption{Dividing shape [6, 17, 128] by bounds [6, 17, 4] gives chunks of shape [1, 1, 32]. Chunk (1, 16, 2) starts at element (1, 16, 64).}
\label{fig:f14-chunkat_3d}
\end{figure}

With the DMA statement, different chunks of data tiled from \texttt{lhs} are moved iteratively and in parallel from the host to the device's \texttt{local} storage. Similarly, the DMA statement manages \texttt{rhs} by moving it in small chunks to \texttt{local} storage for processing.

The declaration \texttt{local s32 [lhs\_\allowbreak{}load.\allowbreak{}span] l1\_\allowbreak{}out;} creates a buffer for each parallel instance, with the same shape as the loaded chunk. The subsequent \texttt{call} invokes the device function \texttt{kernel}, which writes its result to that buffer. A final DMA statement copies the result back to the host output, completing the iteration.

In this code, each parallel thread runs \texttt{17x4} iterations, with each iteration handling a \texttt{1x1x32}-sized chunk of data. The \emph{croqtile program} terminates when all \texttt{6} parallel threads have completed their iterations. It then returns the output buffer to its caller, the host program.


The DMA statement describes data movement and can name its result with a future. Its syntax specifies the operation, source, and destination.


A DMA statement in Croqtile can be an asynchronous entity, meaning that the code following an \textbf{asynchronous DMA} executes in parallel with the DMA statement. For a \emph{asynchronous DMA}, explicit synchronization is required before using the DMA result.

To support these features, the syntax is organized as follows:

\begin{lstlisting}[language=croqtile,caption={Basic Syntax}]
future = dma-op src-expr => dst-expr;
\end{lstlisting}

The statement defines a \textbf{future}-typed variable, which appears on the left side of the operator \texttt{=}. In Croqtile, **a \emph{future} represents both the handle of the asynchronous execution instance and the DMA result**. Therefore, the definition of \emph{future} can be ignored for \emph{synchronized DMA} statements.

The right-hand side of the \textbf{DMA operation} includes the \textbf{operation type} (\texttt{dma-op}), the \textbf{data expression} for the operation source (\texttt{src-expr}), and the \emph{data expression} for the destination (\texttt{dst-expr}). The source and destination are separated by the symbol \texttt{=\textgreater{}}, which indicates the direction of data flow.


\subsection*{Synchronous and Asynchronous}

\textbf{Asynchronous DMA}, or \textbf{Async-DMA}, is also known as \textbf{Non-Blocking DMA} because it does not block the execution of subsequent code. Conversely, \textbf{Synchronous DMA}, or \textbf{Sync-DMA}, is referred to as \textbf{Blocking DMA} because the code does not proceed until the DMA operation is complete. These terms will be used interchangeably in the following text.

In Croqtile, the \texttt{.async} suffix in a \textbf{DMA operation} indicates that the operation is asynchronous. The following code provides an example:

\begin{lstlisting}[language=croqtile,caption={Synchronous and Asynchronous}]
f0 = dma.copy.async data0 => shared; // non-blocking
f1 = dma.copy data1 => shared;       // blocking
...
wait f0;                             // synchronization
\end{lstlisting}

In this example, the DMA operation that produces the future \texttt{f0} is asynchronous, allowing the program to continue executing while the data transfer occurs in the background. The DMA operation that produces \texttt{f1} is synchronous, meaning the program will pause until the data is fully transferred from \texttt{data1} to \texttt{shared} storage. Meanwhile, the data copy of \texttt{data0} might still be in progress.

Executing the statement \texttt{wait f0;} causes the program to check if \texttt{data0} has been fully transferred to \texttt{shared} storage. This synchronizes the main execution flow with the asynchronous DMA of \texttt{f0}. Therefore programmers must make sure there is no "dangling future* left in your code.


\subsection*{Data Expression}

So far, you have only used simple data expressions, either a defined \textbf{spanned data} or a \textbf{storage location}. The \textbf{storage location} is used solely for the \textbf{destination buffer} declaration, as it requires the Croqtile compiler to allocate memory for it. The following code provides an example:

\begin{lstlisting}[language=croqtile,caption={Data Expression}]
global f32 [32, 16] input;
shared f32 [32, 16] output;
dma.copy input => output;  // Copy to a user-declared buffer
dma.copy input => shared;  // Same, but to a compiler-allocated buffer
\end{lstlisting}

In this example, the destination of the second DMA operation is specified as the storage location \texttt{shared}. This requires the compiler to allocate storage accordingly, making it equivalent to the first DMA operation. In practice, you recommend programmers use storage locations as destinations, as Croqtile can deduce the destination shape from the \textbf{DMA statement}.

However, the use of simple \textbf{data expressions} is rare in practice. In the next section, this section introduces the \textbf{ChunkAt} expression to demonstrate how to implement tiling in data expressions.


\subsection*{Future and Wait}

Both \emph{Sync-DMA} and \emph{Async-DMA} operations can be used to define a \emph{future}. As seen in earlier examples, a \emph{future} acts as a handle for asynchronous operations and is typically used as a parameter of the \texttt{wait} statement to synchronize an \emph{Async-DMA}.

Croqtile allows the \texttt{wait} statement to take multiple \emph{future}s. However, using a \emph{future} from a \textbf{Sync-DMA} in a \texttt{wait} statement will result in an error. The following code provides an example:

\begin{lstlisting}[language=croqtile,caption={Future and Wait}]
f0 = dma.copy.async input0 => shared;
f1 = dma.copy.async input1 => shared;
f2 = dma.copy input2 => shared;
...
wait f0, f1;  // Multiple wait
wait f2;      // Error: cannot wait on a sync-dma
\end{lstlisting}

In Croqtile, another use of a \emph{future} is to retrieve the destination buffer of the associated DMA statement. This is also the reason why \textbf{Sync-DMA} is allowed to define a future. Croqtile provides two built-in member functions for \emph{future} variables:

\begin{itemize}[leftmargin=1.6em]
\item \texttt{.span} to retrieve the \textbf{mdspan} of the DMA destination buffer.
\item \texttt{.data} to retrieve the reference of the DMA destination buffer (spanned data).
\end{itemize}

For example:

\begin{lstlisting}[language=croqtile,caption={Future and Wait}]
f0 = dma.copy.async input0 => shared;
f1 = dma.copy input2 => shared;
...
local f32 [f0.span] buffer;
call device_kernel(f1.data, |f1.span|);
\end{lstlisting}

In this example, you use \texttt{.span} built-in member function to retrieve the shape of \texttt{f0}'s destination buffer, and \texttt{.data} to retrieve the \emph{spanned} data of \texttt{f1}'s destination. In addition, the \textbf{ElementCount} operation on the \emph{mdspan} of the destination buffer of \texttt{f1}. The \textbf{ElementCount} operation returns the number of elements in an \textbf{mdspan}, which is useful in many scenarios.


\paragraph*{In short}

You explored both synchronous and asynchronous operations, and saw how Croqtile supports advanced data handling with various \emph{DMA Operation Types}. Finally, you discovered how futures help manage async operations and retrieve buffers, ensuring smooth data flow and synchronization.

In the next section, you will step further to the \emph{Data Expression} to see how to make data tiling/blocking happen.



\section{Coordination: inthreads, Events, and Async Blocks}
\label{sec:2-10}

\guideslotintro{Concurrent regions have to agree on who may proceed. This section introduces \texttt{inthreads}, events and the asynchronous forms through short examples.}

\subsection*{Event arrays for multi-stage pipelines}

Consider a producer that loads tiles and a consumer that computes on them. Before reading a slot, the consumer needs to know its data is ready; before overwriting it, the producer needs to know the previous read has finished. Call these two events \co@full@ and \co@empty@. A pipeline with several reusable slots gives each slot its own pair of events:

\begin{lstlisting}[language=croqtile,caption={Event arrays for multi-stage pipelines}]
shared event full[STAGES], empty[STAGES];
\end{lstlisting}

Each physical buffer slot gets its own \texttt{full}/\texttt{empty} pair. The ring index \texttt{stage = iv\_\allowbreak{}k \% STAGES} maps the unbounded K iteration to a fixed number of physical slots. With four stages, the producer can run several tiles ahead before blocking on \texttt{wait empty}.


\subsection*{Bootstrap protocol}

The consumer must \textbf{trigger} every \texttt{empty} event before the K-loop starts:

\begin{lstlisting}[language=croqtile,caption={Bootstrap protocol}]
foreach {s} in [STAGES] {
  trigger empty[s];
}
\end{lstlisting}

Without this bootstrap, the producer's first \texttt{wait empty[0]} blocks forever --- a deadlock, not a mysterious MMA bug. This is a common pitfall: every bounded buffer protocol needs its events triggered before the first wait.

Section~\ref{sec:3-8-5} develops the full double-buffered and multi-stage pipeline kernels that put these primitives to work. The examples there compose \texttt{inthreads.\allowbreak{}async}, events, \texttt{swap} into complete, runnable matmul pipelines.


\subsection*{Parallel Perspective of Tileflow Programs}

In Croqtile, a \emph{divergent} code region is generated with the keyword \texttt{inthreads}, where a compare expression followed.

Let us have a look at the following example:

\begin{lstlisting}[language=croqtile,caption={Parallel Perspective of Tileflow Programs}]
__co__  void foo() {
  // sequential code
  // start of the SPMD code region
  parallel p by 6 {
    // some uniform code for all threads
    // ...
    inthreads (p <= 2) {
      // divergent code for thread 0, 1, 2 (thread-mask: 000111)
    }
    // threads are forced to be synchronized
    // some uniform code for all threads
    // ...
  }
  // end of the SPMD code region
  // sequential code
}
\end{lstlisting}

You may now re-examine Croqtile tileflow programs from the perspective of parallelism. It consists of

\begin{itemize}[leftmargin=1.6em]
\item \emph{Sequential Code}
\item \emph{SPMD-Styled Parallel Code} enclosed in \texttt{parallel-by} block, which consists of:
\begin{itemize}[leftmargin=1.6em]
\item \emph{Uniform Code}
\item \emph{Divergent Code} enclosed in \texttt{inthreads} block.
\end{itemize}
\end{itemize}

In this example, \texttt{p \textless{}= 2} initiates the \emph{divergent code block}, where from the implementation perspective, it creates a thread mask of \texttt{000111}, which allows only thread \texttt{0}, \texttt{1} and \texttt{2} to execute the code inside. And thread \texttt{3}, \texttt{4} and \texttt{5} will skip this path. \textbf{All threads are forced synchronized after the \texttt{inthreads} block}. In this case, it implies thread \texttt{3}, \texttt{4}, and \texttt{5} waits for the completion of the others before stepping forward.

In Croqtile, the \texttt{inthreads} block looks like a \texttt{if} conditional block in C/C++. However, \texttt{inthreads} explicit claims the code divergence, which is a concept in SPMD program, rather than the conditional block for sequential code. Therefore, the \texttt{inthreads} block:

\begin{itemize}[leftmargin=1.6em]
\item Can only appear in the SPMD code region, which is enclosed in the \texttt{parallel-by} block.
\item Can only have a comparison result related to thread identifiers.
\end{itemize}

Croqtile compiler has the checks for both the comparison expression and \texttt{inthreads}-block placement.


\subsection*{Implicit Masking}

In Croqtile, code between multiple \texttt{parallel-by} levels can be considered as a type of implicit masking. The below code showcases an example:

\begin{lstlisting}[language=croqtile,caption={Implicit Masking}]
__co__  void foo() {
  parallel p by 6 {
    // implicit masking: 'inthreads (q == 0) { ... } '
    parallel q by 2 {
      // code here is for all threads
    }
    // implicit masking: 'inthreads (q == 0) { ... } '
  }
}
\end{lstlisting}

In this example, there exists two levels of parallelization. The code inside \texttt{parallel p by 6} block but outside \texttt{parallel q by 2} blocks is conceptually for the outer parallelization only. This means, only one thread from the sub-level parallelization execute the code, which is equivalent to \texttt{inthreads (q == 0) \{ .\allowbreak{}.\allowbreak{}.\allowbreak{} \}}. Programmers must be aware there is only single thread instance to execute the code part. Only the code inside the inner \texttt{parallel-by} is executed by all threads.


An asynchronous region can advance independently, but a dependency on another region still needs to be expressed. An event marks that dependency: one region triggers it after completing the required work, and another waits before using the result. The event carries readiness information; the data remains in the buffer shared by the regions.


\subsection*{Event Basics}

The meaning of an event varies between programming languages. In Croqtile, an \textbf{event} coordinates asynchronous code blocks at different execution levels.

A Croqtile event has two states, \textbf{SET} and \textbf{UNSET}. A newly declared event starts in the UNSET state. For example:

\begin{lstlisting}[language=croqtile,caption={Declaring a shared event}]
shared event e;
\end{lstlisting}

A \texttt{shared} event is declared inside a \texttt{parallel-by} block and used by its parallel instances. Events can also use \texttt{global} or \texttt{local} storage. The target platform determines the constraints on their declarations and references.

In addition, you may declare an event array following C-style syntax:

\begin{lstlisting}[language=croqtile,caption={Declaring an event array}]
shared event e[4], e1;
\end{lstlisting}

A comma-separated expression can declare both an event and an event array on one line.

The \texttt{trigger} statement changes an event from \textbf{UNSET} to \textbf{SET}:

\begin{lstlisting}[language=croqtile,caption={Triggering an event array}]
shared event e[4], e1;
trigger e, e1;
\end{lstlisting}

The comma-separated statement triggers every event in array \texttt{e} and the single event \texttt{e1}.

The \texttt{wait} statement blocks the current instance while the event is \textbf{UNSET}. Execution resumes once the event becomes \textbf{SET}.

Triggering an event is what unblocks the code waiting on a \emph{wait statement}. For example,

\begin{lstlisting}[language=croqtile,caption={Waiting on an event}]
parallel p by 2 {
  shared event e;                          // init 'e' to 'UNSET'
  inthreads.async (p == 1) { wait e; }     // wait-thread
                                           // once unblocked, change 'e' to 'UNSET'
  inthreads.async (p == 0) { trigger e; }  // trigger-thread
                                           // convert 'e' to 'SET'
  sync.barrier : thread;                   // sync point
}
\end{lstlisting}

In this example, the triggering instance sets \texttt{e}, allowing the waiting instance to continue. Both then reach \texttt{sync.barrier}. The event establishes a dependency between the two asynchronous regions: the trigger must occur before the wait can finish.

In Croqtile semantics, \textbf{an event is always auto-reset by its \texttt{wait} statement}. That means when a \texttt{wait} statement is unblocked, the events that are waited are changed to \emph{UNSET} state.


\paragraph*{In short}

Events let asynchronous \texttt{inthreads} blocks agree on an order. The \texttt{trigger} statements can set events and event arrays to \emph{SET} state, while \texttt{wait} statements block execution, and reset events to \emph{UNSET} until they are triggered.

Programmers must be careful to use Croqtile events, or else it is possible to result in deadlocks, which is hard to debug.



\section{Branching and Early Return}
\label{sec:2-11}

\guideslotintro{Not every branch is known at compile time. This section covers runtime predicates and early return from a parallel region.}

\subsection*{Predicated execution with \texttt{if}}

Croqtile's \texttt{if} behaves like its C counterpart:

\begin{lstlisting}[language=croqtile,caption={Predicated execution with if}]
if (tile_id < total_tiles) {
  // body executes only when the condition is true
}
\end{lstlisting}

All threads in scope evaluate the condition. Threads where it is false skip the body. Threads on the skipped side sit idle while the taken side runs, then vice versa. This is the price of runtime flexibility.

Use \co@if@ for runtime decisions such as bounds checks and conditional accumulation. The condition may depend on the current tile coordinate or input values. Keep it uniform across a cooperating group when the body contains an operation that requires the group to participate together.

Whenever possible, keep \texttt{if} conditions uniform across the execution group: in this example, all 32 threads should agree.


\subsection*{Early return from \texttt{parallel}}

Sometimes you need to early return from a \texttt{parallel} block. The \texttt{yield} keyword generates a \texttt{return;} instruction that leaves the device code for that instance.

\begin{lstlisting}[language=croqtile,caption={Early return from parallel}]
parallel p by BLK_CNT: block {
  parallel q by THR_CNT : thread {
    if (cond) yield;
    ...
  }
}
\end{lstlisting}

\begin{pitfall}[Common pitfall]
The \texttt{yield} keyword currently will \textbf{not} do any compile-time nor runtime checks. Be careful about synchronization safety.

\end{pitfall}



\section{The Four-Step MMA Syntax}
\label{sec:2-12}

\guideslotintro{Contraction is written as four explicit steps. This section reads that syntax once, with the smallest example that shows all four.}

Every contraction kernel follows the same rhythm, and the spelling makes each step visible:

\begin{enumerate}[leftmargin=1.6em]
\item \textbf{\texttt{mma.fill 0.0}}
\item \textbf{\texttt{mma.load}}
\item \textbf{\texttt{mma.\allowbreak{}row.\allowbreak{}row mc, ma, mb}} --- issue the contraction: \textbf{C += A \(\times\) B} into \texttt{mc}.
\item \textbf{\texttt{mma.\allowbreak{}store mc, dst}}
\end{enumerate}

You loop steps 2--3 over K (loading the next K-slice each iteration, accumulating into the same \texttt{mc}), then run step 4 once to flush the completed output tile. You do not declare per-lane layouts; the compiler derives them from the target and your layout choice (\texttt{row.row} here).

The \texttt{.row.row} suffix is a \textbf{layout contract} --- Both operands are row-major. The full set of layout combinations is \texttt{row.row}, \texttt{row.col}, \texttt{col.row}, and \texttt{col.col}; in practice, \texttt{row.row} and \texttt{row.col} cover most kernels. Choosing the wrong variant produces incorrect results.

For two K slices, initialize \co@mc@ once, load the first pair of operands and accumulate, then load the second pair and accumulate into the same \co@mc@. Only then store the output. The load handles describe operand fragments; the accumulator holds the partial result across iterations. This distinction explains why fill and store sit outside the repeated load-and-multiply sequence.



\section{C++ Embedding, Calls, Macros, and the Preprocessor}
\label{sec:2-13}

\guideslotintro{Croqtile lives inside C++: a \file{.co} file can call helpers, inject source text and use the preprocessor. This section covers each mechanism in turn.}

\subsection*{\texttt{\_\_cpp\_\_}: verbatim C++ injection}

The \co@__cpp__@ form inserts a string as C++ source at the point where it appears in a Croqtile function. Use it for a short target operation that needs to sit inside a particular parallel or asynchronous region. The inserted text must be valid in that scope, with all referenced names visible to the target compiler. Croqtile passes the contents through without parsing or rewriting them.

\begin{itemize}[leftmargin=1.6em]
\item \textbf{\texttt{\_\allowbreak{}\_\allowbreak{}cpp\_\allowbreak{}\_\allowbreak{}(".\allowbreak{}.\allowbreak{}.\allowbreak{}")}} --- Ordinary string; best for short one-liners.
\item \textbf{\texttt{\_\allowbreak{}\_\allowbreak{}cpp\_\allowbreak{}\_\allowbreak{}(R"(.\allowbreak{}.\allowbreak{}.\allowbreak{})")}} --- Raw string literal; use for \texttt{asm volatile} where escaping every \texttt{"} would be painful.
\end{itemize}


\subsection*{Macros inside \texttt{\_\_cpp\_\_} strings}

A common mistake: assuming the preprocessor expands macros inside string literals.

\begin{lstlisting}[language=croqtile,caption={A macro inside a cpp string}]
#define PRODUCER_MAXNREG 40
inthreads.async (p1 == 0) {
  __cpp__(R"(asm volatile("setmaxnreg.dec.sync.aligned.u32 PRODUCER_MAXNREG;");)");
}
\end{lstlisting}

This fails --- the preprocessor does not expand inside strings. Use numeric literals inside \texttt{\_\_cpp\_\_} and \texttt{\#if} / \texttt{\#error} outside to enforce consistency:

\begin{lstlisting}[language=croqtile,caption={Its expanded form}]
#define PRODUCER_MAXNREG 40
#if PRODUCER_MAXNREG > 50
#error "Producer maxnreg too high for this tile config"
#endif
inthreads.async (p1 == 0) {
  __cpp__(R"(asm volatile("setmaxnreg.dec.sync.aligned.u32 40;");)");
}
\end{lstlisting}


Croqtile's preprocessor runs before the main compiler pass. Macros expand in both \texttt{\_\_co\_\_} regions and host C++ in the same \texttt{.co} file, so one definition keeps tile geometry and host-side checks aligned.

\begin{table}[htbp]
\centering\small
\begin{tabularx}{\linewidth}{@{}lX@{}}
\toprule
\textbf{Directive} & \textbf{Role} \\
\midrule
\texttt{\#\allowbreak{}define NAME value} & Object-like macro: textual replacement \\
\texttt{\#if} / \texttt{\#elif} / \texttt{\#else} / \texttt{\#endif} & Conditional inclusion \\
\texttt{\#ifdef} / \texttt{\#ifndef} & Shorthand for whether a macro is defined \\
\texttt{\#\allowbreak{}error message} & Force a compile-time failure with a message \\
\bottomrule
\end{tabularx}
\end{table}

Function-like macros (\texttt{\#\allowbreak{}define MAX(a, b) .\allowbreak{}.\allowbreak{}.\allowbreak{}}) are not supported. Use \texttt{constexpr} helpers in ordinary C++.


\subsection*{Tile geometry as macros}

Production matmul sources centralize tile dimensions at the top:

\begin{lstlisting}[language=croqtile,caption={Tile geometry as macros}]
#define MATMUL_WARP_M 64
#define MATMUL_WARP_N 128
#define MATMUL_TILE_K 64
#define MATMUL_WARP_K 16
#define MATMUL_STAGES 4
\end{lstlisting}

The same names appear in \texttt{parallel}, shared-memory declarations, \texttt{foreach} bounds, and host-side verification.


\subsection*{Compile-time assertions with \texttt{\#error}}

\begin{lstlisting}[language=croqtile,caption={Compile-time assertions with error}]
#if MATMUL_SWIZ != (2 * MATMUL_TILE_K)
#error "MATMUL_SWIZ must equal 2 * MATMUL_TILE_K for f16 kernel"
#endif
#if MATMUL_WARP_M != 64
#endif
\end{lstlisting}


\subsection*{Conditional code paths}

\begin{lstlisting}[language=croqtile,caption={Conditional code paths}]
#define PATH0
__co__ foo() {
  #ifdef PATH0
    // path 0 code
  #else
    // path 1 code
  #endif
}
\end{lstlisting}

The preprocessor keeps one branch and discards the other before Croqtile parses the \texttt{\_\_co\_\_} body. Command-line defines: \texttt{croq kernel.\allowbreak{}co -\allowbreak{}DMATMUL\_\allowbreak{}TILE\_\allowbreak{}K=128}.


\subsection*{Calling C++ Functions}

A \co@__co__@ function is not a device function. Calling one from inside another does not produce a device call; the compiler warns that it cannot find a device function of that name. Calls to C++ kernels are the normal case.

\begin{lstlisting}[language=croqtile,caption={Calling C++ Functions}]
void bar() {...}    // C++ kernel function
__co__ void foo() {
  parallel p by 6 {
    call bar();     // Call the C++ function
  }
}
\end{lstlisting}

The example uses the Croqtile \texttt{call} keyword to invoke the existing C++ function \texttt{bar}.


A shared configuration constant often appears in both the host allocation and the kernel declaration. Defining it once as an object-like macro keeps those uses together. The examples below show how replacement text and conditional compilation affect the source before Croqtile checks its shapes and types.


\subsection*{Object-Like Macros}

Croqtile preprocessor supports C/C++ \textbf{object-like macros} to connect both host code and tileflow code. The \emph{object-like macro} performs pure text replacement and accepts no parameters. The following code gives an example:

\begin{lstlisting}[language=croqtile,caption={An object-like macro}]
#define M 256
#define N 32
#define K 64
__co__ auto matmul(f32 [M, N] lhs, f32 [N, K] rhs) { /*...*/ }
void foo() {
  croqtile::f32 a[M][K];
  croqtile::f32 b[N][K];
  // ...
  auto res = matmul(croqtile::make_spanview<2>(a, {M, K}),
                    croqtile::make_spanview<2>(b, {N, K}));
}
\end{lstlisting}

In the code snippet, the inputs of croqtile function \texttt{matmul} are not dynamically shaped. Instead, Croqtile pre-processor substitutes \texttt{M}, \texttt{N}, \texttt{K} with values \texttt{256}, \texttt{32}, \texttt{64} ahead of Croqtile compilation, which produces tileflow function with statically shaped inputs. In this way, programmers make different code consistent and easy to manipulate.

The preprocessor currently supports only macros without parameters, so code like:

\begin{lstlisting}[language={},caption={A macro with parameters}]
#define MAX(a, b) (((a) > (b)) ? (a) : (b))
\end{lstlisting}

would not work.


\subsection*{Comments}

Croqtile support C-style comments, either \texttt{/*...*/} or \texttt{//...}, leveraging the capability croqtile preprocessor has provided.


\subsection*{Conditional Compilation}

Croqtile also support conditional compilation which C/C++ programs used a lot. This includes \texttt{\#if}/\texttt{\#ifdef}/\texttt{\#ifndef}/\texttt{\#else}/\texttt{\#endif}. The below code snippet showcases the usage:

\begin{lstlisting}[language=croqtile,caption={Conditional Compilation}]
#define PATH0
// some host code
__co__ foo() {
#ifdef PATH0
// some code related to PATH0
#else
// other code
#endif
}
// host control
#ifdef PATH0
// ...
#else
// ...
#endif
\end{lstlisting}

In this way, it also make host and croqtile code be controlled within the same preprocess method.

Note that, the capability of croqtile preprocessor is still enhancing. But it is likely that you would not implement full C preprocessing support. Croqtile would not pick up existing C features unless people find it is necessary.


In current implementation of Croqtile functions, it invokes \emph{device function} to implement the specified target computations. This section will learn the detail about the \emph{call statement}s in Croqtile.


\subsection*{Call Device Functions}

To call a device function in the Croqtile function, it requires the call statement explicitly led by \texttt{call}.


\subsection*{Basic Syntax of Call Statements}

The general syntax for a \texttt{call} statement in Croqtile is:

\begin{lstlisting}[language=croqtile,caption={The general call syntax}]
call func-name <optional-template-args> (arguments);
\end{lstlisting}

Here, the keyword \texttt{call} is followed by a function name \texttt{func-name}, an optional \texttt{\textless{}\textgreater{}} enclosed template arguments, which are comma-separated. The \texttt{arguments} are also comma-separated and listed inside \texttt{()} like normal C/C++ functions.

\textbf{Croqtile transpilation process would not apply any check between the callers and the callees}, including the function existence, function signature consistency, and parameter consistency. It delegates such duties to the target compilation process. Programmers must be careful about the conventions between Croqtile function and device function consequently.

In the following example, \texttt{call bar();} in \texttt{mou} is illegal code since it tries to call device function in a code location which is assumed to be host code area.

\begin{lstlisting}[language=croqtile,caption={An illegal call site}]
__co__ void foo() {
  parallel p by 1 {
    call bar();  // ok
  }
}
__co__ void mou() {
  call bar();  // error: not able to call device function from host
}
\end{lstlisting}


\subsection*{The Convention: Allowed Argument Types}

In current implementation, the argument must be either:

\begin{itemize}[leftmargin=1.6em]
\item \emph{Spanned Data} type, or
\item Scalar \emph{Integer} type, or
\item Scalar \emph{Floating-Point} type.
\end{itemize}

Here, the scalar \emph{integer} type only includes \texttt{int}, whereas the scalar \emph{floating-point type} includes \texttt{float}, \texttt{double}, \texttt{half}, \texttt{bfp16}, and \texttt{half8}. However, since the target-platform differs in floating-point support, the \emph{floating-point type} parameter also varies for different target platforms.

In addition, \textbf{operations on floating-point are not supported} by Croqtile till now, considering the floating is not useful for tileflow programs. Instead, they are typically values for device computation, which the \emph{tileflow program}s pass such values directly to device kernels.

The below code showcases an example:

\begin{lstlisting}[language=croqtile,caption={The Convention: Allowed Argument Types}]
__device__ void bar(float *p, int m, int n) {}
__device__ void foo(float *, int, unsigned, double, float) {}
__co__ void foobar(f32 [M, 24] input, int N, float padding) {
  parallel p by 1 {
    shared f32 [14, 7] buffer;
    call bar(input.data, M, buffer.span(0));
    call foo(buffer.data, N, 3, 3.14, padding);
  }
}
\end{lstlisting}

In this example, it passes different data types as arguments from Croqtile function to the device functions, which matches the device parameters exactly. Note that in the device function, the **corresponding parameter type of the \emph{spanned data} argument is simply the pointer of its \_element type\_**, where the shape information is dropped. For example, the \texttt{foo}'s argument \texttt{p} is \texttt{float*}, which corresponds the \emph{spanned data} argument \texttt{input.data}. In Croqtile, you names the parameter as a \textbf{decayed} pointer of the \emph{spanned data}. If using other form like \texttt{float p[]}, it will trigger failure in target compilation stage.

For some types like \texttt{f16}, and \texttt{bf16}, there may not be native target support of such types, it is possible to utilize \texttt{croqtile::f16} and \texttt{croqtile::bf16} to handle such types.


\paragraph*{In short}

A \texttt{call} statement invokes a device function, with or without template arguments. Care is needed since no check is applied at transpilation time, which may result in more confusing error reporting since the check is applied at target compile time with generated functions.



\section{Printing, Assertions, and Compile-Time Inspection}
\label{sec:2-14}

\guideslotintro{The language has its own print and assert forms, at compile time and at runtime. This section shows what each one reports and when to reach for it.}

\subsection*{\texttt{print!} and \texttt{println!}: compile-time inspection}

The bang variants (\texttt{print!}, \texttt{println!}) run at \textbf{compile time}, not at runtime. They print to the compiler's output and are useful for inspecting shapes, extents, and type information without launching a kernel:

\begin{lstlisting}[language=croqtile,caption={A compile-time shape print}]
__co__ void check_shapes(f32 [3, 2] b) {
  print!("shape of b: ");
  println!("b.span = ", b.span);
}
\end{lstlisting}

When you compile this file, the compiler emits:

\begin{lstlisting}[language=croqtile,caption={The printed shape}]
shape of b: b.span = [3, 2]
\end{lstlisting}

String literals are concatenated (\texttt{"wor" "ld!"} becomes \texttt{"world!"}).


\subsection*{\texttt{print} and \texttt{println}: runtime device output}

\begin{lstlisting}[language=croqtile,caption={A runtime inspect kernel}]
__co__ void inspect(s32 [4, 8] data) {
  foreach i in [data.span] {
    println("element ", i, " = ", data.at(i));
  }
}
\end{lstlisting}

Each call takes a comma-separated mix of string literals and expressions. The output order across threads is \textbf{nondeterministic}.

To inspect one execution path, guard the print with a condition on its index:

\begin{lstlisting}[language=croqtile,caption={A print inside a parallel block}]
parallel {px, py} by [8, 16] : block
  parallel {qx, qy} by [16, 16] : thread {
    // ... compute ...
    if (px == 0 && py == 0 && qx == 3 && qy == 5) {
      println("partial sum = ", accum);
    }
  }
\end{lstlisting}

Without the guard you get one line per thread --- thousands of lines for a large grid.


\subsection*{\texttt{assert}: runtime invariant checks}

Croqtile's \texttt{assert} builtin checks an invariant at runtime and aborts the kernel with a message if it fails:

\begin{lstlisting}[language=croqtile,caption={assert: runtime invariant checks}]
assert(stage < MATMUL_STAGES, "stage index out of bounds");
\end{lstlisting}

On the device, this compiles to a \texttt{printf} of the message followed by an abort. Use assertions to catch index-out-of-range, null pointer, and other "this should never happen" conditions early.



\section{Syntax Quick Reference}
\label{sec:2-15}

\guideslotintro{The constructs introduced so far are gathered here as a compact reference, for looking up a spelling or a meaning without re-reading the examples.}

\subsection*{Checklist}

The syntax collected here means the following:

\begin{table}[htbp]
\centering\small
\begin{tabularx}{\linewidth}{@{}lX@{}}
\toprule
\textbf{Syntax} & \textbf{Meaning} \\
\midrule
\texttt{dma.\allowbreak{}copy.\allowbreak{}async .\allowbreak{}.\allowbreak{}.\allowbreak{}} & Start a non-blocking DMA transfer. \\
\texttt{wait f} & Wait until future \texttt{f} is ready to consume. \\
\texttt{future.data} & Access the moved tile. \\
\texttt{future.span} & Access the moved tile's shape. \\
\texttt{dma.\allowbreak{}transp\textless{}.\allowbreak{}.\allowbreak{}.\allowbreak{}\textgreater{}} & Permute dimensions during transfer. \\
\texttt{dma.\allowbreak{}pad\textless{}.\allowbreak{}.\allowbreak{}.\allowbreak{}\textgreater{}} & Pad during transfer with an explicit fill value. \\
\texttt{.zfill} & Zero-fill out-of-bounds positions in partial tiles. \\
\texttt{.\allowbreak{}swiz\textless{}N\textgreater{}} & Choose a shared-memory swizzle layout. \\
\texttt{chunkat(...)} & Select one equal tile from a regular tiling. \\
\texttt{subspan(.\allowbreak{}.\allowbreak{}.\allowbreak{}).\allowbreak{}at(.\allowbreak{}.\allowbreak{}.\allowbreak{})} & Select a tile by explicit tile extents and tile coordinates. \\
\texttt{view(.\allowbreak{}.\allowbreak{}.\allowbreak{}).\allowbreak{}from(.\allowbreak{}.\allowbreak{}.\allowbreak{})} & Select a tile by explicit extents and explicit element offsets. \\
\texttt{tile \# i} & Compose tile and in-tile coordinates into a global coordinate. \\
\texttt{\#tile} & Ask for the extent of the \texttt{tile} axis. \\
\bottomrule
\end{tabularx}
\end{table}

When reading a movement statement, first identify the selected region and its destination, then check when the destination becomes ready. These questions remain the same in a pipeline; the extra buffers simply allow different tiles to be at different points in that sequence.


\subsection*{Summary of New Syntax}

\begin{table}[htbp]
\centering\small
\begin{tabularx}{\linewidth}{@{}lX@{}}
\toprule
\textbf{Syntax} & \textbf{Meaning} \\
\midrule
\texttt{parallel p by N} & N-way parallelism, index \texttt{p} runs 0..N-1 concurrently \\
\texttt{parallel \{px, py\} by [M, N]} & Multi-dimensional parallel level   (\texttt{p} and \texttt{q} correspond to one parallel level) \\
\texttt{parallel p by A, q by B} & Comma-separated parallel levels   (\texttt{p} and \texttt{q} correspond to different parallel levels) \\
\texttt{parallel q by N :\allowbreak{} thread} & Map to threads within a block \\
\texttt{foreach index = \{m, n, k\} in [a, b, c]} & Named tuple destructuring in a \texttt{foreach} \\
\texttt{lhs.span(0)} & Extract one dimension of a tensor's shape \\
\texttt{p \# m} & Compose outer tile index \texttt{p} with inner offset \texttt{m} \\
\bottomrule
\end{tabularx}
\end{table}

Section 3.7 replaces scalar arithmetic in the inner loop with matrix contraction operations.


\paragraph*{Syntax at a glance}

\begin{table}[htbp]
\centering\small
\begin{tabularx}{\linewidth}{@{}lX@{}}
\toprule
\textbf{Syntax} & \textbf{Meaning} \\
\midrule
\texttt{if (expr) \{ .\allowbreak{}.\allowbreak{}.\allowbreak{} \}} & Predicated execution --- runtime conditional, divergent threads masked \\
\texttt{inthreads.\allowbreak{}async (condition)} & Structured concurrent region --- compile-time thread partitioning \\
\texttt{shared event name[N]} & Declare N synchronization tokens \\
\texttt{trigger name} & Record that a condition is met \\
\texttt{wait name} & Block until the corresponding \texttt{trigger} fires \\
\texttt{tile\_\allowbreak{}id = tile\_\allowbreak{}iter \#\allowbreak{} block\_\allowbreak{}id} & Compose indices for tile striping \\
\texttt{int total\_\allowbreak{}tiles = expr} & Local integer variable \\
\texttt{yield} & Early return from \texttt{parallel} block \\
\bottomrule
\end{tabularx}
\end{table}


\paragraph*{Syntax at a glance}

\begin{table}[htbp]
\centering\small
\begin{tabularx}{\linewidth}{@{}lX@{}}
\toprule
\textbf{Syntax} & \textbf{Meaning} \\
\midrule
\texttt{dma.\allowbreak{}copy.\allowbreak{}async src =\textgreater{} dst} & Non-blocking copy (returns immediately) \\
\texttt{dma.any} & Placeholder future (no transfer in flight yet) \\
\texttt{swap(f0, f1)} & Exchange two future handles without copying data \\
\texttt{rotate(f0, f1, f2)} & Cycle three future handles \\
\texttt{with tile\_\allowbreak{}k in N \{ .\allowbreak{}.\allowbreak{}.\allowbreak{} \}} & Scoped tile axis binding with extent N \\
\texttt{foreach tile\_\allowbreak{}k(1:\allowbreak{})} & Iterate starting from index 1 \\
\texttt{\_\allowbreak{}\_\allowbreak{}co\_\allowbreak{}\_\allowbreak{} auto fn(.\allowbreak{}.\allowbreak{}.\allowbreak{})} & Return type inferred from \texttt{return} statement \\
\bottomrule
\end{tabularx}
\end{table}



\guidechapternext{You can now read most \file{.co} files. Chapter~\ref{ch:programming-guide} explains why the constructs are shaped this way and covers the advanced forms; Appendix~\ref{app:quickref} is a compact lookup table.}

